\documentclass[10pt,a4paper]{article}
% Оформление: ivanshemidko/algebra-notes, s2/main.tex,
% commit d52e89ab46b905aac035e556794af91aec7c86ef.
\usepackage{fontspec}
% Project-local installation keeps local and CI font files identical.
\setmainfont{cmunrm.ttf}[
  Path=fonts/,
  BoldFont=cmunbx.ttf,
  ItalicFont=cmunti.ttf,
  BoldItalicFont=cmunbi.ttf
]
\usepackage[english,russian]{babel}
\usepackage{amsmath,amssymb,mathtools,amsthm}
% Unicode maps for traditional Computer Modern math fonts (search/copy text).
\pdfvariable gentounicode=1
\protected\def\pdfglyphtounicode{\pdfextension glyphtounicode}
\input{glyphtounicode.tex}
% The standard map omits CMEX's sized delimiter glyph names.
\ExplSyntaxOn
\clist_map_inline:nn {big,Big,bigg,Bigg} {
  \pdfglyphtounicode{parenleft#1}{0028}
  \pdfglyphtounicode{parenright#1}{0029}
  \pdfglyphtounicode{braceleft#1}{007B}
  \pdfglyphtounicode{braceright#1}{007D}
  \pdfglyphtounicode{bracketleft#1}{005B}
  \pdfglyphtounicode{bracketright#1}{005D}
  \pdfglyphtounicode{floorleft#1}{230A}
  \pdfglyphtounicode{floorright#1}{230B}
  \pdfglyphtounicode{ceilingleft#1}{2308}
  \pdfglyphtounicode{ceilingright#1}{2309}
  \pdfglyphtounicode{angbracketleft#1}{27E8}
  \pdfglyphtounicode{angbracketright#1}{27E9}
  \pdfglyphtounicode{slash#1}{002F}
  \pdfglyphtounicode{backslash#1}{005C}
}
\ExplSyntaxOff
\usepackage{geometry}
\geometry{left=22mm,right=22mm,top=25mm,bottom=27mm,heightrounded}
\usepackage{enumitem}
\usepackage[most]{tcolorbox}
\tcbuselibrary{theorems,breakable,skins}
\usepackage{graphicx,caption}
\captionsetup{font=small,labelfont=bf,hypcap=false}

\definecolor{cbBlueBack}{RGB}{246,249,252}
\definecolor{cbGreenBack}{RGB}{246,251,248}
\definecolor{cbOrangeBack}{RGB}{252,249,242}
\definecolor{cbPurpleBack}{RGB}{249,247,252}
\definecolor{cbGrayBack}{RGB}{249,249,249}
\definecolor{cbBlue}{RGB}{42,87,132}
\definecolor{cbGreen}{RGB}{52,110,72}
\definecolor{cbOrange}{RGB}{140,98,0}
\definecolor{cbPurple}{RGB}{93,78,140}
\definecolor{cbGray}{RGB}{80,80,80}

\newtheoremstyle{mystatement}{8pt}{8pt}{\itshape}{}{\bfseries}{.}{\newline}{}
\newtheoremstyle{mydefinition}{8pt}{8pt}{\normalfont}{}{\bfseries}{.}{0.5em}{}
\theoremstyle{mystatement}
\newtheorem{theorem}{Теорема}[part]
\newtheorem{lemma}{Лемма}[part]
\newtheorem{cor}{Следствие}[part]
\theoremstyle{mydefinition}
\newtheorem{definition}{Определение}[part]
\newtheorem{example}{Пример}[part]
\numberwithin{equation}{part}
\counterwithin{figure}{part}
\tcbset{
  enhanced,sharp corners,boxrule=0.5pt,
  left=9pt,right=9pt,top=6pt,bottom=6pt,
  before skip=6pt,after skip=6pt,breakable
}
\tcolorboxenvironment{theorem}{colback=cbBlueBack,colframe=cbGray!55,borderline west={1.6pt}{0pt}{cbBlue}}
\tcolorboxenvironment{lemma}{colback=cbBlueBack,colframe=cbGray!55,borderline west={1.6pt}{0pt}{cbBlue}}
\tcolorboxenvironment{cor}{colback=cbPurpleBack,colframe=cbGray!55,borderline west={1.6pt}{0pt}{cbPurple}}
\tcolorboxenvironment{definition}{colback=cbGreenBack,colframe=cbGreen!55!cbGray,borderline west={1.6pt}{0pt}{cbGreen}}
\tcolorboxenvironment{example}{colback=cbOrangeBack,colframe=cbOrange!45!cbGray,borderline west={1.6pt}{0pt}{cbOrange}}

\usepackage{hyperref}
\hypersetup{
  colorlinks=true,linkcolor=cbBlue,citecolor=cbBlue,urlcolor=cbBlue,
  pdftitle={Дискретная математика · Коллоквиум 2},
  pdfauthor={},pdfsubject={},pdfcreator={LuaLaTeX}
}
\linespread{1.0}
\setlength{\parindent}{0pt}
\setlength{\parskip}{4pt}
\allowdisplaybreaks[1]
\providecommand{\tightlist}{\setlength{\itemsep}{0pt}\setlength{\parskip}{0pt}}

\usepackage{titlesec}
\titleformat{\part}[block]{\Large\bfseries\centering}{}{0pt}{}
\titlespacing*{\part}{0pt}{0pt}{12pt}
\counterwithin{section}{part}
\renewcommand{\thepart}{\arabic{part}}
\usepackage{tocloft}
\setlength{\cftsecnumwidth}{2.8em}
\setlength{\cftsubsecnumwidth}{3.6em}
\setcounter{tocdepth}{2}
\usepackage{fancyhdr}
\pagestyle{fancy}
\fancyhf{}
\fancyhead[L]{\nouppercase{\leftmark}}
\fancyfoot[C]{\thepage}
\let\oldpart\part
\renewcommand{\part}[1]{\oldpart{#1}\markboth{#1}{}}
\renewcommand{\sectionmark}[1]{}
\usepackage{etoolbox}
\pretocmd{\part}{\newpage}{}{}

\usepackage{tikz}
\usetikzlibrary{arrows.meta,positioning,calc,patterns,fit,shapes.geometric}
% Общий язык учебных схем; штрихи, двойной контур и метки читаются без цвета.
\tikzset{
  gi vertex/.style={circle,draw=black,fill=white,minimum size=5mm,inner sep=1pt},
  gi small/.style={circle,draw=black,fill=white,minimum size=2.5mm,inner sep=0pt},
  gi selected/.style={gi vertex,double,fill=black!12},
  gi edge/.style={draw=black,line width=.65pt},
  gi emphasis/.style={draw=cbBlue,line width=1.25pt},
  gi alternate/.style={draw=black,densely dashed,line width=1pt},
  gi removed/.style={draw=black!55,densely dotted,line width=.85pt},
  gi arrow/.style={gi edge,-{Stealth[length=2mm]}},
  gi flow/.style={-{Stealth[length=2mm]},thick},
  gi group/.style={draw=black!60,dashed,rounded corners,inner sep=6pt},
  gi note/.style={align=center,font=\small},
  gi box/.style={draw,rounded corners,align=center,inner sep=5pt,font=\small},
  gi match/.style={draw=cbBlue,line width=2.2pt},
  gi free/.style={circle,draw=black,fill=black!12,minimum size=5mm,inner sep=1pt,double},
  gi c1/.style={draw=cbBlue,line width=1pt},
  gi c2/.style={draw=cbOrange,densely dashed,line width=1pt},
  gi c3/.style={draw=cbGreen,densely dotted,line width=1.3pt},
  gi interval/.style={line width=1.4pt,{|[width=2.5mm]}-{|[width=2.5mm]}}
}
% Неплавающая композиция: рисунок и подпись переходят на страницу вместе.
\newenvironment{graphillustration}[2]{%
  \par\addvspace{6pt}\noindent\begin{minipage}{\linewidth}%
  \def\gicaption{#1}\def\gilabel{#2}\centering\normalfont\small
}{%
  \par\captionof{figure}{\gicaption}\label{\gilabel}%
  \end{minipage}\par\addvspace{6pt}%
}
\DeclareMathOperator{\comp}{c}
\DeclareMathOperator{\odd}{odd}
\DeclareMathOperator{\dfc}{def}
\DeclareMathOperator{\diag}{diag}
\DeclareMathOperator{\cl}{cl}
\DeclareMathOperator{\rank}{rank}
\newcommand{\N}{\mathbb N}
\newcommand{\Z}{\mathbb Z}
\newcommand{\R}{\mathbb R}


% Билеты используют счётчик part, но компактные строки оглавления без двойной нумерации.
\newcommand{\ticket}[1]{%
  \clearpage\refstepcounter{part}%
  \markboth{Билет \thepart. #1}{}%
  \addcontentsline{toc}{part}{Билет \thepart. #1}%
  {\centering\Large\bfseries Билет \thepart. #1\par}%
  \vspace{12pt}%
}
\makeatletter
\renewcommand*{\l@part}[2]{%
  \addpenalty{-100}\addvspace{5pt}%
  \@dottedtocline{0}{0em}{0em}{\bfseries #1}{#2}%
}
\makeatother

\begin{document}
\thispagestyle{empty}
\begin{center}
{\LARGE\bfseries Дискретная математика · III семестр}\par\medskip
{\Large Теория графов. Коллоквиум 2}\par\medskip
{\large Паросочетания · Теория Рамсея · Совершенные графы}
\end{center}

\clearpage
\tableofcontents
\clearpage
\phantomsection
\section*{Обозначения и соглашения}
\markboth{Обозначения и соглашения}{}
\addcontentsline{toc}{section}{Обозначения и соглашения}
\begingroup
\setlength{\parskip}{4pt}
\titlespacing*{\subsection}{0pt}{9pt}{4pt}
\tcbset{before skip=4pt,after skip=4pt,top=4pt,bottom=4pt}
\setlist[description]{font=\normalfont\bfseries,leftmargin=0pt,
  labelindent=0pt,style=unboxed,itemsep=3pt,parsep=0pt}
Все графы конечные и простые (без петель и кратных рёбер), если не оговорено иное;
\(G=(V,E)\), \(n=|V|\), \(m=|E|\).

\subsection*{Основные обозначения}
{\renewcommand{\arraystretch}{1.22}
\begin{tabular}{@{}p{0.30\linewidth}@{\hspace{1em}}p{\dimexpr0.70\linewidth-1em\relax}@{}}
\textbf{Обозначение} & \textbf{Значение} \\
\hline
\(N(v)\), \(N(A)\) & Соседи вершины; \(N(A)=\bigcup_{a\in A}N(a)\) \\
\(G-S\), \(G[S]\) & Удаление множества вершин; индуцированный подграф \\
\(H\le G\) & \(H\) — индуцированный подграф \(G\) \\
\(\overline G\) & Дополнение: рёбра ровно между несмежными в \(G\) вершинами \\
\(V(M)\) & Множество концов рёбер \(M\subseteq E\) \\
\(M\mathbin{\triangle}M'\) & Симметрическая разность (в лекциях \(M\oplus M'\)) \\
\(\odd(G)\), \(\comp(G)\) & Число нечётных компонент; число всех компонент связности \\
\(K_n\), \(K_{p,q}\), \(C_n\), \(P_n\) & Полный, полный двудольный граф, цикл и путь на \(n\) вершинах
\end{tabular}\par}

\subsection*{Множества вершин}
\begin{definition}[Независимое множество]
\(I\subseteq V\) \emph{независимо}, если никакие две его вершины не смежны.
\(\alpha(G)=\max|I|\) — \emph{число независимости}.
\end{definition}
\begin{definition}[Вершинное покрытие]
\(C\subseteq V\) — \emph{вершинное покрытие}, если у каждого ребра хотя бы один конец в \(C\).
\(\beta(G)=\min|C|\) (в лекциях \(\mathrm{vc}(G)\)).
\end{definition}
\begin{definition}[Клика]
\(K\subseteq V\) — \emph{клика}, если \(G[K]\) полный. \(\omega(G)=\max|K|\) — \emph{кликовое число};
\(\chi(G)\) — хроматическое число. Заметим \(\alpha(G)=\omega(\overline G)\).
\end{definition}

\subsection*{Множества рёбер}
\begin{definition}[Паросочетание]
\(M\subseteq E\) — \emph{паросочетание}, если никакие два его ребра не имеют общего конца.
\(\alpha'(G)=\max|M|\) (в лекциях \(\mathrm{mm}(G)\)).
\(M\) \emph{покрывает} вершину \(v\), если \(v\in V(M)\); остальные вершины \emph{свободны}.
\end{definition}
\begin{definition}[Рёберное покрытие]
\(F\subseteq E\) — \emph{рёберное покрытие}, если каждая вершина — конец ребра из \(F\).
\(\beta'(G)=\min|F|\); существует лишь при \(\delta(G)>0\).
\end{definition}
\begin{definition}[Совершенное и почти совершенное паросочетание]
Паросочетание \emph{совершенно}, если покрывает все вершины (тогда \(|M|=n/2\)),
и \emph{почти совершенно}, если покрывает все вершины, кроме одной.
\end{definition}
\begin{definition}[Чередующийся и дополняющий путь]
Простой путь \emph{\(M\)-чередующийся}, если его рёбра поочерёдно лежат и не лежат в \(M\).
\emph{\(M\)-дополняющий} путь — \(M\)-чередующийся путь положительной длины с двумя
свободными концами; рёбер не из \(M\) в нём на одно больше.
\end{definition}
\begin{tcolorbox}[colback=cbOrangeBack,colframe=cbOrange!45!cbGray,
  borderline west={1.6pt}{0pt}{cbOrange}]
\textbf{Соглашение.} «Максимальное» паросочетание (независимое множество) — наибольшее
по размеру. Не расширяемое добавлением элемента называем \emph{максимальным по
включению}; оно может быть меньше (рис.~\ref{fig:notation-matchings}).
\end{tcolorbox}
\begin{graphillustration}{Путь $P_4$: среднее ребро — паросочетание, максимальное по включению (ни одно ребро не добавить), но не наибольшее. Два крайних ребра — совершенное паросочетание.}{fig:notation-matchings}
\begin{tikzpicture}[font=\small]
\begin{scope}
\node[gi free] (a) at (0,0) {$a$};
\node[gi vertex] (b) at (1.3,0) {$b$};
\node[gi vertex] (c) at (2.6,0) {$c$};
\node[gi free] (d) at (3.9,0) {$d$};
\draw[gi edge] (a) -- (b);
\draw[gi match] (b) -- (c);
\draw[gi edge] (c) -- (d);
\node[gi note] at (1.95,-0.8) {$M=\{bc\}$: максимально по включению,\\$a,b,c,d$ — дополняющий путь};
\end{scope}
\begin{scope}[shift={(6.5,0)}]
\node[gi vertex] (a) at (0,0) {$a$};
\node[gi vertex] (b) at (1.3,0) {$b$};
\node[gi vertex] (c) at (2.6,0) {$c$};
\node[gi vertex] (d) at (3.9,0) {$d$};
\draw[gi match] (a) -- (b);
\draw[gi edge] (b) -- (c);
\draw[gi match] (c) -- (d);
\node[gi note] at (1.95,-0.8) {$M'=\{ab,cd\}$: совершенное,\\$|M'|=\alpha'(P_4)=2$};
\end{scope}
\end{tikzpicture}
\end{graphillustration}


\subsection*{Раскраски рёбер}
\begin{definition}[Рёберная раскраска]
Раскраска рёбер \(K_t\) в \(k\) цветов — функция \(\varphi\colon E(K_t)\to\{1,\ldots,k\}\).
При \(k=2\) она равносильна графу \(G=(V,\varphi^{-1}(1))\): рёбра цвета 2 образуют
\(\overline G\), клики цвета 2 — независимые множества \(G\).
\end{definition}
\textbf{На рисунках:} рёбра паросочетания — толстые синие; свободные вершины — двойной
серый контур; цвета 1, 2, 3 — сплошная синяя, штриховая оранжевая, пунктирная зелёная линии.

\subsection*{Исходные факты}
\begin{lemma}[О рукопожатиях]
\(\sum_{v}d(v)=2m\); число вершин нечётной степени чётно.
\end{lemma}
\begin{lemma}[Рёбра и компоненты]
В любом графе \(m\ge n-\comp(G)\).
\end{lemma}
\begin{lemma}[Клика и раскраска]
\(\chi(G)\ge\omega(G)\).
\end{lemma}

\endgroup
\ticket{Теорема Бержа о дополняющих цепях}
\label{ticket:1}
\begin{lemma}[Симметрическая разность паросочетаний]\label{lem:symdiff}
Пусть \(M_1,M_2\) — паросочетания графа \(G\), \(H=(V,\,M_1\mathbin{\triangle}M_2)\).
Тогда каждая компонента \(H\) — изолированная вершина, простой путь или
простой цикл чётной длины, и вдоль каждого пути и цикла рёбра \(M_1\) и \(M_2\)
чередуются.
\end{lemma}
\begin{proof}
Каждая вершина инцидентна не более чем одному ребру \(M_1\) и не более чем
одному ребру \(M_2\), поэтому \(d_H(v)\le2\), и если \(d_H(v)=2\), то два ребра
при \(v\) из разных паросочетаний. Граф со степенями \(\le2\) — объединение
путей и циклов. Чередование означает, что цикл имеет чётную длину.
\end{proof}
\begin{graphillustration}{Компоненты $M\triangle M'$: чётный цикл и пути. Толстые рёбра — $M$, штриховые — $M'$. В пути справа рёбер $M'$ больше; это $M$-дополняющий путь.}{fig:01-symdiff}
\begin{tikzpicture}[font=\small]
\begin{scope}
\node[gi vertex] (c1) at (0,1) {};
\node[gi vertex] (c2) at (1.2,1) {};
\node[gi vertex] (c3) at (1.8,0) {};
\node[gi vertex] (c4) at (1.2,-1) {};
\node[gi vertex] (c5) at (0,-1) {};
\node[gi vertex] (c6) at (-0.6,0) {};
\draw[gi match] (c1) -- (c2); \draw[gi alternate] (c2) -- (c3);
\draw[gi match] (c3) -- (c4); \draw[gi alternate] (c4) -- (c5);
\draw[gi match] (c5) -- (c6); \draw[gi alternate] (c6) -- (c1);
\node[gi note] at (0.6,-1.7) {чётный цикл: поровну};
\end{scope}
\begin{scope}[shift={(3.6,0)}]
\node[gi vertex] (p1) at (0,0.6) {};
\node[gi vertex] (p2) at (1.1,0.6) {};
\node[gi vertex] (p3) at (2.2,0.6) {};
\draw[gi match] (p1) -- (p2); \draw[gi alternate] (p2) -- (p3);
\node[gi note] at (1.1,-0.3) {чётный путь:\\поровну};
\end{scope}
\begin{scope}[shift={(7.2,0)}]
\node[gi free] (q1) at (0,0.6) {};
\node[gi vertex] (q2) at (1.1,0.6) {};
\node[gi vertex] (q3) at (2.2,0.6) {};
\node[gi free] (q4) at (3.3,0.6) {};
\draw[gi alternate] (q1) -- (q2); \draw[gi match] (q2) -- (q3);
\draw[gi alternate] (q3) -- (q4);
\node[gi note] at (1.65,-0.3) {нечётный путь: $M'$ на 1 больше,\\концы свободны в $M$};
\end{scope}
\end{tikzpicture}
\end{graphillustration}


\begin{theorem}[Берж, 1957]
Паросочетание \(M\) максимально (\(|M|=\alpha'(G)\)) тогда и только тогда,
когда в \(G\) нет \(M\)-дополняющих путей.
\end{theorem}
\begin{proof}
\(\Rightarrow\). Пусть есть \(M\)-дополняющий путь \(P=a_1a_2\ldots a_{2k}\).
Рёбра \(a_2a_3,\ldots,a_{2k-2}a_{2k-1}\) лежат в \(M\), а
\(a_1a_2,a_3a_4,\ldots,a_{2k-1}a_{2k}\) — нет. «Инвертируем» путь:
\[
M'=M\mathbin{\triangle}E(P)=\bigl(M\setminus E(P)\bigr)\cup\bigl(E(P)\setminus M\bigr).
\]
Внутренние вершины пути по-прежнему покрыты ровно одним ребром, концы
\(a_1,a_{2k}\) были свободны и стали покрыты, остальные вершины не затронуты.
Значит, \(M'\) — паросочетание и \(|M'|=|M|+1\). Противоречие с максимальностью.

\(\Leftarrow\). Пусть \(M\)-дополняющих путей нет, а \(M'\) — максимальное
паросочетание. По лемме~\ref{lem:symdiff} граф \(H=(V,M\mathbin{\triangle}M')\) —
объединение чётных циклов и путей с чередующимися рёбрами. В цикле и в пути
чётной длины рёбер \(M\) и \(M'\) поровну. Если бы \(|M'|>|M|\), нашлась бы
компонента-путь, где рёбер \(M'\) больше: путь нечётной длины, начинающийся и
кончающийся рёбрами \(M'\setminus M\). Его концы не покрыты \(M\): конец пути
имеет степень 1 в \(H\), и если бы он был покрыт ребром \(e\in M\), то
\(e\notin M'\) (у вершины уже есть ребро \(M'\)), т.\,е. \(e\in H\) и степень была бы 2.
Получен \(M\)-дополняющий путь — противоречие. Следовательно, \(|M|=|M'|\).
\end{proof}
\begin{graphillustration}{Инверсия дополняющего пути $a_1\ldots a_6$: рёбра пути меняют статус, свободные концы становятся покрытыми, $|M|$ растёт на 1.}{fig:01-augment}
\begin{tikzpicture}[font=\small]
\foreach \i/\x in {1/0,2/1.1,3/2.2,4/3.3,5/4.4,6/5.5} {
  \ifnum\i=1 \node[gi free] (a\i) at (\x,0) {$a_{\i}$};
  \else\ifnum\i=6 \node[gi free] (a\i) at (\x,0) {$a_{\i}$};
  \else \node[gi vertex] (a\i) at (\x,0) {$a_{\i}$}; \fi\fi
}
\draw[gi edge] (a1) -- (a2); \draw[gi match] (a2) -- (a3);
\draw[gi edge] (a3) -- (a4); \draw[gi match] (a4) -- (a5);
\draw[gi edge] (a5) -- (a6);
\node[gi note] at (2.75,-0.8) {$M$: 2 ребра на пути};
\node[gi note] at (6.6,0) {$\Longrightarrow$};
\begin{scope}[shift={(7.6,0)}]
\foreach \i/\x in {1/0,2/1.1,3/2.2,4/3.3,5/4.4,6/5.5} \node[gi vertex] (b\i) at (\x,0) {$a_{\i}$};
\draw[gi match] (b1) -- (b2); \draw[gi edge] (b2) -- (b3);
\draw[gi match] (b3) -- (b4); \draw[gi edge] (b4) -- (b5);
\draw[gi match] (b5) -- (b6);
\node[gi note] at (2.75,-0.8) {$M\triangle E(P)$: 3 ребра};
\end{scope}
\end{tikzpicture}
\end{graphillustration}


\begin{cor}[Алгоритмическая идея]
Начиная с пустого паросочетания и увеличивая его вдоль дополняющих путей,
пока они есть, получаем максимальное паросочетание: процесс заканчивается
не позже чем через \(n/2\) шагов, а по теореме Бержа остановка возможна
только на максимальном паросочетании.
\end{cor}

\begin{example}
На рис.~\ref{fig:notation-matchings} паросочетание \(\{bc\}\) максимально по
включению, но путь \(a,b,c,d\) дополняющий; инверсия даёт \(\{ab,cd\}\).
Именно поэтому жадное добавление рёбер не гарантирует максимальность.
\end{example}

\ticket{Теорема Кёнига: паросочетание и вершинное покрытие в двудольном графе}
\label{ticket:2}
Пусть \(G=(V_1,V_2,E)\) — двудольный граф с долями \(V_1,V_2\).
\begin{lemma}[Слабая двойственность]
В любом графе \(\alpha'(G)\le\beta(G)\).
\end{lemma}
\begin{proof}
Пусть \(M\) — паросочетание, \(C\) — вершинное покрытие. Каждое ребро \(M\)
имеет конец в \(C\), а одна вершина \(C\) покрывает не более одного ребра \(M\)
(рёбра \(M\) не имеют общих концов). Значит, \(|M|\le|C|\).
\end{proof}
Для треугольника \(\alpha'(K_3)=1<2=\beta(K_3)\), так что в общем случае
равенства нет. В двудольном графе оно есть.

\begin{theorem}[Кёниг, 1931]
Если \(G\) двудолен, то \(\alpha'(G)=\beta(G)\): размер максимального
паросочетания равен размеру минимального вершинного покрытия.
\end{theorem}
\begin{proof}
По лемме достаточно построить вершинное покрытие \(B\) размера \(|M|\),
где \(M\) — максимальное паросочетание.

Пусть \(U_1\subseteq V_1\) — свободные вершины доли \(V_1\), \(U_2\subseteq V_2\) —
свободные вершины доли \(V_2\). Назовём вершину \emph{достижимой}, если до неё
есть \(M\)-чередующийся путь из \(U_1\), начинающийся ребром не из \(M\)
(вершины \(U_1\) достижимы путём длины 0). Положим
\[
Y_1=\text{достижимые из }V_1\setminus U_1,\quad Z_1=V_1\setminus(U_1\cup Y_1),\qquad
Y_2=\text{достижимые из }V_2,\quad Z_2=V_2\setminus Y_2 .
\]
\textit{Свойства.}
\begin{enumerate}\tightlist
\item \(Y_2\cap U_2=\varnothing\): иначе путь до свободной \(y\in Y_2\) был бы
  \(M\)-дополняющим, что невозможно по теореме Бержа. Значит, \(U_2\subseteq Z_2\).
\item Путь из \(U_1\) переходит из \(V_1\) в \(V_2\) по рёбрам не из \(M\), а из
  \(V_2\) в \(V_1\) — по рёбрам \(M\). Поэтому партнёр по \(M\) каждой
  \(y\in Y_2\) лежит в \(Y_1\) (путь продолжается ребром \(M\)), и каждая
  вершина \(Y_1\) получена так. Итак, \(M\) — биекция \(Y_1\leftrightarrow Y_2\);
  следовательно, партнёры вершин \(Z_1\) (они покрыты по определению) лежат в \(Z_2\).
\item Если \(x\in U_1\cup Y_1\) и \(xy\in E\setminus M\), то \(y\in Y_2\): путь до \(x\)
  кончается ребром \(M\) (или пуст), и его можно продлить ребром \(xy\).
\end{enumerate}
Положим \(B=Z_1\cup Y_2\).

\emph{\(B\) — покрытие.} Ребро \(xy\) (\(x\in V_1\), \(y\in V_2\)) с \(x\in Z_1\)
покрыто. Если \(x\in U_1\cup Y_1\): при \(xy\notin M\) имеем \(y\in Y_2\) по
свойству 3; при \(xy\in M\) вершина \(x\in Y_1\) и её партнёр \(y\in Y_2\) по
свойству 2.

\emph{\(|B|=|M|\).} Каждое ребро \(M\) соединяет либо \(Y_1\) с \(Y_2\), либо
\(Z_1\) с \(Z_2\) (свойство 2), т.\,е. имеет ровно один конец в \(B\);
и каждая вершина \(B\) покрыта (свойство 1 и определение \(Z_1\)).
Значит, \(|B|=|M|=\alpha'(G)\), и \(\beta(G)\le\alpha'(G)\).
\end{proof}
\begin{graphillustration}{Структура из доказательства Кёнига. Толстые рёбра — $M$, крестиком отмечены невозможные рёбра. Покрытие $B=Z_1\cup Y_2$ выделено двойным контуром групп.}{fig:02-konig}
\begin{tikzpicture}[font=\small]
\node[gi free] (u1) at (0,2) {};
\node[gi vertex] (y1a) at (2.2,2) {};
\node[gi vertex] (y1b) at (3.2,2) {};
\node[gi vertex] (z1a) at (5.6,2) {};
\node[gi vertex] (z1b) at (6.6,2) {};
\node[gi free] (u2) at (0,0) {};
\node[gi vertex] (y2a) at (2.2,0) {};
\node[gi vertex] (y2b) at (3.2,0) {};
\node[gi vertex] (z2a) at (5.6,0) {};
\node[gi vertex] (z2b) at (6.6,0) {};
\draw[gi match] (y1a) -- (y2a); \draw[gi match] (y1b) -- (y2b);
\draw[gi match] (z1a) -- (z2a); \draw[gi match] (z1b) -- (z2b);
\draw[gi edge] (u1) -- (y2a); \draw[gi edge] (y1a) -- (y2b);
\draw[gi edge] (z1a) -- (y2b); \draw[gi edge] (z1b) -- (z2a);
\draw[gi removed] (u1) -- (u2) node[pos=.5,fill=white,inner sep=1pt] {$\times$};
\draw[gi removed] (y1b) -- (z2a) node[pos=.5,fill=white,inner sep=1pt] {$\times$};
\node[gi group,fit=(u1),label=above:$U_1$] {};
\node[gi group,fit=(y1a)(y1b),label=above:$Y_1$] {};
\node[gi group,double,fit=(z1a)(z1b),label=above:{$Z_1\subseteq B$}] {};
\node[gi group,fit=(u2),label=below:$U_2$] {};
\node[gi group,double,fit=(y2a)(y2b),label=below:{$Y_2\subseteq B$}] {};
\node[gi group,fit=(z2a)(z2b),label=below:$Z_2$] {};
\node[gi note,anchor=west] at (7.6,2) {$V_1$};
\node[gi note,anchor=west] at (7.6,0) {$V_2$};
\node[gi note,anchor=west] at (8.4,1) {$U_1$—$U_2$: дополняющий путь;\\$U_1\cup Y_1$—$Z_2$: $Z_2$ стала бы\\достижимой.};
\end{tikzpicture}
\end{graphillustration}


\begin{cor}[Двудольный граф, \(\delta>0\)]
\(\alpha(G)=\beta'(G)\).
\end{cor}
\begin{proof}
По билету~\ref{ticket:3} \(\alpha+\beta=n=\alpha'+\beta'\), а \(\beta=\alpha'\) по Кёнигу.
\end{proof}
\begin{example}
В \(C_6\) (двудолен) \(\alpha'=3\) и покрытие из трёх вершин одной доли.
В \(C_5\) (не двудолен) \(\alpha'=2\), а \(\beta=3\): две вершины покрывают
не более четырёх рёбер из пяти.
\end{example}

\ticket{Вершинные покрытия, независимые множества и рёберные покрытия}
\label{ticket:3}
\begin{lemma}[Дополнение независимого множества]
\(I\subseteq V\) независимо тогда и только тогда, когда \(V\setminus I\) —
вершинное покрытие.
\end{lemma}
\begin{proof}
\(\Rightarrow\). Если ребро \(uv\) не покрыто \(V\setminus I\), то \(u,v\in I\) —
два смежных элемента независимого множества, противоречие.
\(\Leftarrow\). Если \(u,v\in I\) смежны, то ребро \(uv\) не имеет концов в
\(V\setminus I\), т.\,е. не покрыто, противоречие.
\end{proof}
\begin{cor}
\(\alpha(G)+\beta(G)=n\).
\end{cor}
\begin{proof}
Отображение \(I\mapsto V\setminus I\) — биекция между независимыми множествами
и покрытиями, \(|V\setminus I|=n-|I|\). Наибольшему \(I\) соответствует
наименьшее покрытие.
\end{proof}
\begin{graphillustration}{Пример: независимое множество $I=\{a,c\}$ (двойной контур) и его дополнение — вершинное покрытие $\{b,d,e\}$. Каждое ребро имеет конец вне $I$.}{fig:03-cover}
\begin{tikzpicture}[font=\small]
\node[gi selected] (a) at (90:1.2) {$a$};
\node[gi vertex] (b) at (162:1.2) {$b$};
\node[gi selected] (c) at (234:1.2) {$c$};
\node[gi vertex] (d) at (306:1.2) {$d$};
\node[gi vertex] (e) at (18:1.2) {$e$};
\draw[gi edge] (a) -- (b) -- (c) -- (d) -- (e) -- (a);
\node[gi note,anchor=west] at (2.2,0.3) {$C_5$: $\alpha=2$, $\beta=3$,\\$\alpha+\beta=5=n$};
\end{tikzpicture}
\end{graphillustration}


\begin{theorem}[Галлаи, 1959]
Если \(\delta(G)>0\) (нет изолированных вершин), то \(\alpha'(G)+\beta'(G)=n\).
\end{theorem}
\begin{proof}
\(\le\). Пусть \(M\) — максимальное паросочетание, \(|M|=\alpha'\), \(U=V\setminus V(M)\)
— свободные вершины, \(|U|=n-2\alpha'\). Так как \(\delta>0\), для каждой
\(u\in U\) выберем инцидентное ей ребро; пусть \(F\) — множество выбранных рёбер,
\(|F|\le|U|\). Тогда \(M\cup F\) — рёберное покрытие:
\[
\beta'\le|M\cup F|\le|M|+|U|=\alpha'+n-2\alpha'=n-\alpha'.
\]
\(\ge\). Пусть \(L\) — минимальное рёберное покрытие, \(|L|=\beta'\),
\(H=(V,L)\). В \(H\) нет изолированных вершин, поэтому в каждой компоненте
\(H\) есть ребро; выберем по одному ребру из каждой компоненты. Рёбра из разных
компонент не имеют общих концов, так что получилось паросочетание \(N\) графа
\(G\) размера \(\comp(H)\). Отсюда, используя \(e(H)\ge n-\comp(H)\),
\[
\alpha'\ge|N|=\comp(H),\qquad \beta'=|L|=e(H)\ge n-\comp(H)\ge n-\alpha'.\qedhere
\]
\end{proof}
\begin{graphillustration}{Обе половины доказательства Галлаи. Слева: максимальное паросочетание $M$ (толстые) и по ребру (штрих) для каждой свободной вершины; $|M\cup F|=2+2=n-\alpha'$. Справа: минимальное рёберное покрытие — лес звёзд; по ребру из каждой компоненты (толстые) дают паросочетание.}{fig:03-gallai}
\begin{tikzpicture}[font=\small]
\begin{scope}
\node[gi vertex] (b) at (1.2,0.5) {$b$};
\node[gi vertex] (a) at (0,0.5) {$a$};
\node[gi free] (c) at (1.2,1.6) {$c$};
\node[gi free] (u) at (1.2,-0.6) {$u$};
\node[gi vertex] (d) at (2.6,0.5) {$d$};
\node[gi vertex] (e) at (3.8,0.5) {$e$};
\draw[gi match] (a) -- (b); \draw[gi match] (d) -- (e);
\draw[gi edge] (b) -- (d);
\draw[gi alternate] (c) -- (b); \draw[gi alternate] (u) -- (b);
\node[gi note] at (1.9,-1.4) {$M\cup F$, $F=\{cb,ub\}$};
\end{scope}
\begin{scope}[shift={(6.6,0)}]
\node[gi vertex] (s) at (0,0.5) {};
\node[gi vertex] (s1) at (-0.8,1.2) {};
\node[gi vertex] (s2) at (-0.8,-0.2) {};
\node[gi vertex] (s3) at (0.9,1.2) {};
\draw[gi match] (s) -- (s1); \draw[gi edge] (s) -- (s2); \draw[gi edge] (s) -- (s3);
\node[gi vertex] (t) at (2.0,0.5) {};
\node[gi vertex] (t1) at (3.0,0.5) {};
\draw[gi match] (t) -- (t1);
\node[gi vertex] (r) at (4.2,1.0) {};
\node[gi vertex] (r1) at (4.2,0) {};
\node[gi vertex] (r2) at (5.1,0.5) {};
\draw[gi match] (r) -- (r2); \draw[gi edge] (r1) -- (r2);
\node[gi note] at (2.2,-1.4) {$H=(V,L)$: $\comp(H)=3$, $e(H)\ge n-3$};
\end{scope}
\end{tikzpicture}
\end{graphillustration}

\begin{tcolorbox}[colback=cbGrayBack,colframe=cbGray!55]
\textbf{Сводка.} Всегда \(\alpha+\beta=n\), \(\alpha'\le\beta\);
при \(\delta>0\) также \(\alpha'+\beta'=n\) и \(\alpha\le\beta'\)
(каждое ребро покрытия покрывает не более одной вершины независимого множества).
Для двудольных графов неравенства превращаются в равенства (Кёниг).
\end{tcolorbox}
\begin{example}
Для звезды \(K_{1,3}\): \(\alpha=3,\ \beta=1,\ \alpha'=1,\ \beta'=3\); обе суммы равны \(4\).
Для \(C_5\): \(\alpha=2,\ \beta=3,\ \alpha'=2,\ \beta'=3\).
\end{example}

\ticket{Теорема Татта о совершенном паросочетании}
\label{ticket:4}
\begin{definition}[Условие Татта]
\(\odd(G)\) — число нечётных компонент связности графа \(G\).
Граф удовлетворяет \emph{условию Татта}, если \(\odd(G-S)\le|S|\) для всех
\(S\subseteq V\). Множество \(S\) с \(\odd(G-S)>|S|\) называется
\emph{множеством Татта}.
\end{definition}
\begin{theorem}[Татт, 1947]
В графе \(G\) есть совершенное паросочетание тогда и только тогда, когда
\(\odd(G-S)\le|S|\) для всех \(S\subseteq V\).
\end{theorem}
\textit{Необходимость.} Пусть \(M\) — совершенное паросочетание и \(S\subseteq V\).
Нечётную компоненту \(Q\) графа \(G-S\) нельзя разбить на пары внутри себя,
значит, некоторая её вершина сопоставлена ребром \(M\) вершине вне \(Q\);
другие компоненты \(G-S\) с \(Q\) не смежны, поэтому партнёр лежит в \(S\).
Партнёры для разных компонент различны, так что \(\odd(G-S)\le|S|\).
\begin{graphillustration}{Множество Татта: после удаления $S=\{s\}$ остаются три нечётные компоненты, а партнёров в $S$ хватит только одной. Совершенного паросочетания нет.}{fig:04-tutte-set}
\begin{tikzpicture}[font=\small]
\node[gi selected] (s) at (0,0) {$s$};
\node[gi vertex] (a) at (-2.2,-1.4) {};
\node[gi vertex] (b1) at (0,-1.4) {};
\node[gi vertex] (b2) at (-0.6,-2.4) {};
\node[gi vertex] (b3) at (0.6,-2.4) {};
\node[gi vertex] (c) at (2.2,-1.4) {};
\draw[gi edge] (s) -- (a); \draw[gi edge] (s) -- (b1); \draw[gi edge] (s) -- (c);
\draw[gi edge] (b1) -- (b2) -- (b3) -- (b1);
\node[gi group,fit=(a)] {};
\node[gi group,fit=(b1)(b2)(b3)] {};
\node[gi group,fit=(c)] {};
\node[gi note,anchor=west] at (3.2,-1.2) {$\odd(G-S)=3>1=|S|$};
\end{tikzpicture}
\end{graphillustration}


\textit{Достаточность} (доказательство Ловаса). Две вспомогательные леммы.
\begin{lemma}\label{lem:odd-edge}
Если \(uv\notin E\), то \(\odd(G+uv)\le\odd(G)\).
\end{lemma}
\begin{proof}
Если \(u,v\) в одной компоненте, компоненты не меняются. Иначе две компоненты
сливаются: чёт\(+\)чёт \(=\) чёт, нечёт\(+\)чёт \(=\) нечёт (число нечётных
не меняется), нечёт\(+\)нечёт \(=\) чёт (уменьшается на 2).
\end{proof}
\begin{lemma}
Если \(G\) удовлетворяет условию Татта, то и \(G+uv\) удовлетворяет.
\end{lemma}
\begin{proof}
Для \(S\subseteq V\): если \(u\in S\) или \(v\in S\), то \((G+uv)-S=G-S\);
иначе \((G+uv)-S=(G-S)+uv\), и по лемме~\ref{lem:odd-edge}
\(\odd\bigl((G+uv)-S\bigr)\le\odd(G-S)\le|S|\).
\end{proof}
\begin{proof}[Доказательство достаточности]
Предположим, что теорема неверна. Среди контрпримеров (условие Татта
выполнено, совершенного паросочетания нет) на данном множестве вершин \(V\)
выберем граф \(G\) с \emph{максимальным} числом рёбер. Тогда
\begin{itemize}\tightlist
\item \(n\) чётно: при \(S=\varnothing\) имеем \(\odd(G)\le0\);
\item для любого \(e\notin E\) граф \(G+e\) удовлетворяет условию Татта и имеет
  больше рёбер, значит, \emph{в \(G+e\) есть совершенное паросочетание};
\item \(G\ne K_n\) (в \(K_n\) с чётным \(n\) совершенное паросочетание есть).
\end{itemize}
Пусть \(U=\{u\in V: d(u)=n-1\}\) — вершины, смежные со всеми; \(U\ne V\).

\begin{lemma}\label{lem:cliques}
\(G-U\) — дизъюнктное объединение клик.
\end{lemma}
\textit{Доказательство леммы.} Иначе в одной компоненте \(G-U\) найдутся \(x,y,z\) с \(xy,yz\in E\), \(xz\notin E\)
(возьмём кратчайший путь между несмежными вершинами компоненты). Так как
\(y\notin U\), есть \(w\ne y\) с \(yw\notin E\). По выбору \(G\) существуют совершенные
паросочетания \(M_1\) графа \(G+xz\) и \(M_2\) графа \(G+yw\); поскольку в \(G\)
совершенного паросочетания нет, \(xz\in M_1\) и \(yw\in M_2\).
Граф \(H=(V,M_1\mathbin{\triangle}M_2)\) состоит из изолированных вершин и
чётных чередующихся циклов (лемма билета~\ref{ticket:1}; путей нет: оба
паросочетания совершенны), причём \(xz\) и \(yw\) лежат на циклах.

\emph{Случай 1: \(xz\) и \(yw\) на разных циклах \(C_1\), \(C_2\).} Возьмём рёбра
\(M_2\) на \(C_1\), рёбра \(M_1\) на \(C_2\) и рёбра \(M_1\) на остальных вершинах.
Это совершенное паросочетание без \(xz\) и \(yw\), т.\,е. в \(G\) — противоречие.

\emph{Случай 2: на одном цикле \(C\).} Пойдём по \(C\) от \(y\) через ребро \(yw\);
не умаляя общности (\(x\) и \(z\) симметричны: обе смежны с \(y\)), раньше
встретим \(z\), так что \(C=y\,w\ldots z\,x\ldots y\). Путь
\[
P=\underbrace{x\ldots y}_{\text{дуга }C}\;\; yz\;\;\underbrace{z\ldots w}_{\text{дуга }C}
\]
проходит все вершины \(C\) и не использует \(xz\), \(yw\); его рёбра лежат в \(G\).
Число вершин \(C\) чётно, поэтому рёбра \(P\) через одно дают совершенное
паросочетание \(V(C)\). Добавив \(M_1\) вне \(C\), получим совершенное паросочетание
\(G\) — противоречие. Лемма доказана.
\begin{graphillustration}{Случай 2 леммы о кликах. Слева цикл $C$ из $M_1\triangle M_2$: толстые рёбра — $M_1$, штриховые — $M_2$; $xz\in M_1$ и $yw\in M_2$ — добавленные рёбра. Справа путь $P$: дуга $x\ldots y$, ребро $yz\in E$, дуга $z\ldots w$; рёбра $P$ через одно (толстые) дают совершенное паросочетание на $V(C)$.}{fig:04-cycle}
\begin{tikzpicture}[font=\small]
\begin{scope}
\foreach \i/\a/\l in {0/90/y,1/30/w,2/-30/p,3/-90/z,4/-150/x,5/150/q}
  \node[gi vertex] (c\i) at (\a:1.4) {$\l$};
\draw[gi alternate] (c0) -- node[above right,inner sep=1pt] {\scriptsize новое} (c1);
\draw[gi match] (c1) -- (c2);
\draw[gi alternate] (c2) -- (c3);
\draw[gi match,dash pattern=on 4pt off 2pt] (c3) -- node[below,inner sep=2pt] {\scriptsize новое} (c4);
\draw[gi alternate] (c4) -- (c5);
\draw[gi match] (c5) -- (c0);
\draw[gi edge,densely dotted] (c0) -- (c3);
\node[gi note] at (0,-2.2) {$C$; пунктир $yz\in E(G)$};
\end{scope}
\node[gi note] at (3.0,0) {$\Longrightarrow$};
\begin{scope}[shift={(6,0)}]
\foreach \i/\a/\l in {0/90/y,1/30/w,2/-30/p,3/-90/z,4/-150/x,5/150/q}
  \node[gi vertex] (d\i) at (\a:1.4) {$\l$};
\draw[gi match] (d4) -- (d5);
\draw[gi edge] (d5) -- (d0);
\draw[gi match] (d0) -- (d3);
\draw[gi edge] (d3) -- (d2);
\draw[gi match] (d2) -- (d1);
\draw[gi removed] (d0) -- (d1);
\draw[gi removed] (d3) -- (d4);
\node[gi note] at (0,-2.2) {$P=x\,q\,y\,z\,p\,w$};
\end{scope}
\end{tikzpicture}
\end{graphillustration}


\emph{Завершение.} По условию \(\odd(G-U)\le|U|\). В каждой чётной клике \(G-U\)
разобьём вершины на пары. В каждой нечётной клике разобьём на пары все вершины,
кроме одной, а её сопоставим своей вершине \(U\) (вершины \(U\) смежны со всеми, а
различных вершин хватает). Осталось чётное число вершин \(U\) (всего \(n\) чётно,
а покрыто уже чётное число); они попарно смежны, разобьём их на пары.
Получено совершенное паросочетание \(G\) — противоречие.
\end{proof}
\begin{graphillustration}{Завершение доказательства Татта: $G-U$ — несвязные клики. Чётные клики разбиваются на пары внутри; в каждой нечётной одна вершина уходит к своей вершине $U$; остаток $U$ — клика с чётным числом вершин.}{fig:04-finish}
\begin{tikzpicture}[font=\small]
\node[gi vertex] (u1) at (0,2) {};
\node[gi vertex] (u2) at (1.2,2) {};
\node[gi vertex] (u3) at (2.4,2) {};
\node[gi vertex] (u4) at (3.6,2) {};
\node[gi group,fit=(u1)(u2)(u3)(u4),label=above:{$U$ (смежны со всеми)}] {};
\draw[gi match] (u3) -- (u4);
\node[gi vertex] (a1) at (-0.6,0) {};
\node[gi vertex] (a2) at (0.4,0) {};
\node[gi vertex] (a3) at (-0.1,-0.8) {};
\draw[gi edge] (a1) -- (a2) -- (a3) -- (a1);
\draw[gi match] (a2) -- (a3);
\node[gi vertex] (b1) at (1.8,0) {};
\draw[gi match] (a1) -- (u1);
\draw[gi match] (b1) -- (u2);
\node[gi vertex] (c1) at (3.0,0) {};
\node[gi vertex] (c2) at (4.0,0) {};
\node[gi vertex] (c3) at (3.0,-0.9) {};
\node[gi vertex] (c4) at (4.0,-0.9) {};
\draw[gi edge] (c1) -- (c2) -- (c4) -- (c3) -- (c1) -- (c4); \draw[gi edge] (c2) -- (c3);
\draw[gi match] (c1) -- (c2); \draw[gi match] (c3) -- (c4);
\node[gi group,fit=(a1)(a2)(a3)] {};
\node[gi group,fit=(b1)] {};
\node[gi group,fit=(c1)(c2)(c3)(c4)] {};
\node[gi note,anchor=west] at (4.8,-0.4) {клики $G-U$:\\две нечётные, одна чётная};
\end{tikzpicture}
\end{graphillustration}


\begin{cor}[Петерсен]
Любой 3-регулярный граф без мостов имеет совершенное паросочетание.
\end{cor}
\begin{proof}[Набросок (по книге)]
Пусть \(S\subseteq V\), \(Q\) — нечётная компонента \(G-S\). Сумма степеней в \(Q\)
равна \(3|Q|\) — нечётна, а внутренние рёбра дают чётный вклад, поэтому
из \(Q\) в \(S\) ведёт нечётное число рёбер, т.\,е. не меньше \(3\) (одно ребро было бы
мостом). Итого \(3\odd(G-S)\le\) рёбер, входящих в \(S\), \(\le3|S|\).
\end{proof}

\ticket{Формула Бержа для дефицита графа}
\label{ticket:5}
\begin{definition}[Дефицит]
\(\dfc(G)=n-2\alpha'(G)\) — число вершин, не покрытых максимальным
паросочетанием. Эквивалентно \(\alpha'(G)=\dfrac{n-\dfc(G)}{2}\).
\(\dfc(G)=0\) тогда и только тогда, когда есть совершенное паросочетание.
\end{definition}
\begin{theorem}[Формула Бержа, 1958]
\[
\dfc(G)=\max_{S\subseteq V}\bigl(\odd(G-S)-|S|\bigr),\qquad
\alpha'(G)=\frac12\Bigl(n-\max_{S\subseteq V}\bigl(\odd(G-S)-|S|\bigr)\Bigr).
\]
\end{theorem}
При \(S=\varnothing\) выражение равно \(\odd(G)\ge0\), так что максимум неотрицателен.
Теорема Татта — частный случай: \(\dfc(G)=0\iff\odd(G-S)\le|S|\) для всех \(S\).
\begin{proof}
\(\ge\). Пусть \(M\) — максимальное паросочетание, \(S\subseteq V\),
\(Q_1,\ldots,Q_t\) — нечётные компоненты \(G-S\), \(t=\odd(G-S)\). В каждой \(Q_i\)
внутренние рёбра \(M\) не покрывают хотя бы одну вершину \(u_i\); она либо свободна,
либо сопоставлена вершине из \(S\). Сопоставленных не больше \(|S|\), значит,
свободных среди \(u_1,\ldots,u_t\) не меньше \(t-|S|\): \(\dfc(G)\ge\odd(G-S)-|S|\).

\(\le\). Пусть \(k=\max_S(\odd(G-S)-|S|)\). При \(k=0\) по теореме Татта есть
совершенное паросочетание и \(\dfc(G)=0\). Пусть \(k>0\).

\emph{Чётность.} Для любого \(S\): \(n-|S|=\sum|V(Q)|\) по компонентам \(G-S\),
и по модулю 2 это \(\odd(G-S)\). Значит, \(\odd(G-S)-|S|\equiv n\pmod2\),
в частности \(k\equiv n\pmod 2\).

Построим \(H\): добавим к \(G\) клику \(W\) из \(k\) новых вершин и соединим каждую
вершину \(W\) со всеми вершинами \(G\). Тогда \(v(H)=n+k\) чётно. Проверим условие
Татта для \(H\) и \(T\subseteq V(H)\):
\begin{enumerate}\tightlist
\item \(T=\varnothing\): \(H\) связен и чётен, \(\odd(H)=0\).
\item \(T\ne\varnothing\), \(W\not\subseteq T\): в \(H-T\) осталась вершина \(W\), смежная
  со всеми, поэтому \(H-T\) связен и \(\odd(H-T)\le1\le|T|\).
\item \(W\subseteq T\), \(T=W\cup S\), \(S\subseteq V\): \(H-T=G-S\) и
  \(\odd(H-T)=\odd(G-S)\le k+|S|=|T|\) по определению \(k\).
\end{enumerate}
По теореме Татта в \(H\) есть совершенное паросочетание \(N\), \(|N|=(n+k)/2\).
Выбросим из \(N\) рёбра, инцидентные \(W\) — их не больше \(k\). Останется
паросочетание \(G\) размера \(\ge\frac{n+k}{2}-k=\frac{n-k}{2}\), т.\,е.
\(\dfc(G)\le n-(n-k)=k\).
\end{proof}
\begin{graphillustration}{Граф $H$ из доказательства (здесь $n=6$, $k=2$): новые вершины $W$ образуют клику и смежны со всеми вершинами $G$. Совершенное паросочетание $H$ использует не более $k$ рёбер при $W$; остальные его рёбра — паросочетание $G$.}{fig:05-construction}
\begin{tikzpicture}[font=\small]
\node[draw,ellipse,minimum width=4.2cm,minimum height=2.2cm] (G) at (0,0) {};
\node[gi note] at (0,-1.4) {$G$};
\foreach \i/\x/\y in {1/-1.3/0.4,2/-0.3/0.6,3/0.9/0.5,4/1.2/-0.4,5/-1.1/-0.5,6/0.0/-0.5}
  \node[gi small] (g\i) at (\x,\y) {};
\node[gi vertex] (w1) at (4.2,0.8) {$w_1$};
\node[gi vertex] (w2) at (4.2,-0.8) {$w_2$};
\draw[gi edge] (w1) -- (w2);
\node[gi group,fit=(w1)(w2),label=right:{$W$, $|W|=k$}] {};
\foreach \i in {1,...,6} {\draw[gi edge,black!30] (w1) -- (g\i); \draw[gi edge,black!30] (w2) -- (g\i);}
\draw[gi match] (w1) -- (g3); \draw[gi match] (w2) -- (g4);
\draw[gi match] (g1) -- (g2); \draw[gi match] (g5) -- (g6);
\end{tikzpicture}
\end{graphillustration}


\begin{example}
Для звезды \(K_{1,3}\) с центром \(c\): \(S=\{c\}\) даёт \(\odd(G-S)-|S|=3-1=2\),
поэтому \(\dfc=2\), \(\alpha'=(4-2)/2=1\). Для графа на рис.~\ref{fig:04-tutte-set}:
\(n=6\), \(\dfc\ge3-1=2\), и действительно \(\alpha'=2\).
\end{example}
\begin{tcolorbox}[colback=cbGrayBack,colframe=cbGray!55]
\textbf{Как пользоваться.} Чтобы доказать \(\alpha'(G)\le r\), достаточно предъявить
\(S\) с \(\odd(G-S)-|S|\ge n-2r\) («свидетель»); чтобы доказать \(\alpha'(G)\ge r\) —
предъявить паросочетание. Формула гарантирует, что свидетель всегда существует.
\end{tcolorbox}

\ticket{Декомпозиция Эдмондса–Галлаи}
\label{ticket:6}
\subsection*{Фактор-критические графы}
\begin{definition}
Граф \(G\) \emph{фактор-критический}, если для любой вершины \(u\) граф \(G-u\)
имеет совершенное паросочетание. Тогда \(n\) нечётно, \(\dfc(G)=1\) и
\(\alpha'(G-u)=\alpha'(G)=\frac{n-1}2\) для всех \(u\).
Примеры: \(K_1\), нечётные циклы, \(K_{2k+1}\).
\end{definition}
\begin{lemma}[Галлаи]\label{lem:gallai}
Если \(G\) связен и \(\alpha'(G-u)=\alpha'(G)\) для всех \(u\in V\), то \(G\) фактор-критический.
\end{lemma}
\begin{proof}
По формуле Бержа выберем \(S\) с \(\odd(G-S)-|S|=\dfc(G)\) и предположим \(S\ne\varnothing\).
Пусть \(v\in S\), \(S'=S\setminus\{v\}\). Тогда \((G-v)-S'=G-S\) и
\[
\dfc(G-v)\ge\odd(G-S)-|S'|=\dfc(G)+1 .
\]
Подставляя \(\dfc=v(\cdot)-2\alpha'\): \((n-1)-2\alpha'(G-v)\ge n-2\alpha'(G)+1\), т.\,е.
\(\alpha'(G-v)\le\alpha'(G)-1\) — противоречие с условием. Значит, \(S=\varnothing\) и
\(\dfc(G)=\odd(G)\le1\) (граф связен). Совершенного паросочетания нет (его
удаление любой вершины уменьшило бы \(\alpha'\)), поэтому \(\dfc(G)=1\),
\(n\) нечётно и \(\alpha'(G-u)=\alpha'(G)=\frac{n-1}{2}\): у \(G-u\) совершенное паросочетание.
\end{proof}

\subsection*{Множества \(D\), \(A\), \(C\)}
\begin{definition}[Разбиение Галлаи–Эдмондса]
Для графа \(G\) положим:
\begin{itemize}\tightlist
\item \(D(G)\) — вершины \(u\), для которых существует максимальное паросочетание,
  не покрывающее \(u\) («иногда свободные»);
\item \(A(G)\) — вершины вне \(D(G)\), смежные хотя бы с одной вершиной \(D(G)\);
\item \(C(G)=V\setminus(A(G)\cup D(G))\).
\end{itemize}
\end{definition}
Вершины \(A\cup C\) покрыты \emph{любым} максимальным паросочетанием.
\begin{lemma}[Stability lemma]\label{lem:stability}
Если \(a\in A(G)\), то \(D(G-a)=D(G)\), \(A(G-a)=A(G)\setminus\{a\}\), \(C(G-a)=C(G)\)
и \(\alpha'(G-a)=\alpha'(G)-1\).
\end{lemma}
\begin{proof}
Так как \(a\notin D(G)\), каждое максимальное паросочетание \(M\) покрывает \(a\),
и \(M\) без ребра при \(a\) — паросочетание \(G-a\); при этом \(\alpha'(G-a)\ne\alpha'(G)\)
(иначе максимальное паросочетание \(G-a\) было бы максимальным в \(G\) и не покрывало \(a\)).
Значит, \(\alpha'(G-a)=\alpha'(G)-1\). Равенства для \(A\) и \(C\) следуют из
\(D(G-a)=D(G)\) (по определению \(A,C\) через \(D\)). Докажем его.

\(\supseteq\). Пусть \(v\in D(G)\), \(M_v\) — максимальное паросочетание \(G\), не
покрывающее \(v\), \(ax\in M_v\). Тогда \(M_v\setminus\{ax\}\) имеет размер
\(\alpha'(G)-1=\alpha'(G-a)\), т.\,е. максимально в \(G-a\), и не покрывает \(v\).

\(\subseteq\). Пусть \(M'\) — максимальное паросочетание \(G-a\), не покрывающее
\(v\notin D(G)\). Возьмём \(w\in D(G)\), смежную с \(a\), и максимальное паросочетание
\(M_w\) графа \(G\), не покрывающее \(w\); оно покрывает \(v\) (так как \(v\notin D(G)\)).
В \(M_w\mathbin{\triangle}M'\) компонента \(P\), содержащая \(v\), — путь с началом \(v\),
первое ребро из \(M_w\). Вершина \(a\) не покрыта \(M'\), поэтому лежит не более
чем на одном ребре разности и может быть только концом \(P\).
\begin{enumerate}\tightlist
\item[(а)] \(P\) кончается ребром \(M'\): тогда \(M_w\mathbin{\triangle}E(P)\) — паросочетание
  того же размера, не покрывающее \(v\), т.\,е. \(v\in D(G)\) — противоречие.
\item[(б)] \(P\) кончается ребром \(M_w\) в вершине \(a\): тогда
  \((M_w\mathbin{\triangle}E(P))\cup\{aw\}\) — максимальное паросочетание \(G\),
  не покрывающее \(v\) (вершина \(w\) была свободна) — противоречие.
\item[(в)] \(P\) кончается ребром \(M_w\), \(a\notin P\): тогда \(P\) — \(M'\)-дополняющий путь в
  \(G-a\) (последняя вершина не покрыта \(M'\), так как степень 1), и
  \(|M'\mathbin{\triangle}E(P)|=|M'|+1\) — противоречие с максимальностью \(M'\).\qedhere
\end{enumerate}
\end{proof}
\begin{graphillustration}{Три случая в доказательстве $D(G-a)\subseteq D(G)$. Путь $P$ из $M_w\triangle M'$ начинается в $v$; толстые рёбра — $M_w$, штриховые — $M'$. Вершина $w$ свободна в $M_w$, вершина $a$ — в $M'$.}{fig:06-stability}
\begin{tikzpicture}[font=\small]
\begin{scope}
\node[gi vertex] (v) at (0,0) {$v$};
\node[gi vertex] (p1) at (1,0) {};
\node[gi vertex] (p2) at (2,0) {};
\node[gi vertex] (p3) at (3,0) {};
\draw[gi match] (v) -- (p1); \draw[gi alternate] (p1) -- (p2); \draw[gi match] (p2) -- (p3);
\node[gi vertex] (p4) at (4,0) {};
\draw[gi alternate] (p3) -- (p4);
\node[gi note] at (2,-0.8) {(а) конец ребром $M'$};
\end{scope}
\begin{scope}[shift={(5.2,0)}]
\node[gi vertex] (v) at (0,0) {$v$};
\node[gi vertex] (p1) at (1,0) {};
\node[gi vertex] (p2) at (2,0) {};
\node[gi vertex] (a) at (3,0) {$a$};
\node[gi free] (w) at (3,1.1) {$w$};
\draw[gi match] (v) -- (p1); \draw[gi alternate] (p1) -- (p2); \draw[gi match] (p2) -- (a);
\draw[gi edge,densely dotted] (a) -- (w);
\node[gi note] at (1.6,-0.8) {(б) конец $a$: добавим $aw$};
\end{scope}
\begin{scope}[shift={(9.4,0)}]
\node[gi vertex] (v) at (0,0) {$v$};
\node[gi vertex] (p1) at (1,0) {};
\node[gi vertex] (p2) at (2,0) {};
\node[gi vertex] (p3) at (3,0) {};
\draw[gi match] (v) -- (p1); \draw[gi alternate] (p1) -- (p2); \draw[gi match] (p2) -- (p3);
\node[gi note] at (1.5,-0.8) {(в) $M'$-дополняющий путь};
\end{scope}
\end{tikzpicture}
\end{graphillustration}


\subsection*{Теорема о структуре}
\begin{theorem}[Галлаи 1964, Эдмондс 1965]
Пусть \(D_1,\ldots,D_t\) — компоненты связности \(G[D(G)]\), \(A=A(G)\), \(C=C(G)\). Тогда
\begin{enumerate}\tightlist
\item \(G[C]\) имеет совершенное паросочетание;
\item каждый граф \(D_i\) фактор-критический;
\item любое максимальное паросочетание \(M\) состоит из совершенного паросочетания
  \(G[C]\), почти совершенных паросочетаний всех \(D_i\) и рёбер, соединяющих каждую
  вершину \(A\) с вершинами \emph{различных} компонент \(D_i\);
\item \(\dfc(G)=t-|A|\), \(\ \alpha'(G)=\dfrac{n-t+|A|}{2}\).
\end{enumerate}
\end{theorem}
\begin{proof}
Удаляя вершины \(A\) по одной и применяя лемму~\ref{lem:stability}, получаем
\[
D(G-A)=D(G),\quad A(G-A)=\varnothing,\quad C(G-A)=C,\quad\alpha'(G-A)=\alpha'(G)-|A|.
\]
Так как \(A(G-A)=\varnothing\), в \(G-A\) нет рёбер между \(D\) и \(C\); значит,
\(D_1,\ldots,D_t\) и \(G[C]\) — объединение компонент \(G-A\).

1) Вершины \(C=C(G-A)\) покрыты любым максимальным паросочетанием \(G-A\), а
рёбер из \(C\) в \(D\) нет. Значит, его часть внутри \(C\) — совершенное паросочетание \(G[C]\).

2) Для \(u\in D_i\) есть максимальное паросочетание \(M_u\) графа \(G-A\), не
покрывающее \(u\); его часть внутри компоненты \(D_i\) — максимальное паросочетание
\(D_i\), не покрывающее \(u\). Итак, \(\alpha'(D_i-u)=\alpha'(D_i)\) для всех \(u\), и по
лемме Галлаи~\ref{lem:gallai} \(D_i\) фактор-критический.

3) Пусть \(M\) максимально в \(G\), \(M'\) — \(M\) без рёбер, инцидентных \(A\). Тогда
\(|M'|\ge|M|-|A|=\alpha'(G-A)\), т.\,е. \(M'\) максимально в \(G-A\), и все неравенства —
равенства: каждая вершина \(A\) покрыта своим ребром, и ни одно ребро \(M\) не
соединяет две вершины \(A\). По пунктам 1–2 \(M'\) состоит из совершенного
паросочетания \(G[C]\) и почти совершенных паросочетаний \(D_i\) (в фактор-критическом
графе максимальное паросочетание почти совершенно). Свободных вершин в каждой
\(D_i\) ровно одна, а в \(C\) их нет, поэтому рёбра из \(A\) идут к различным \(D_i\).

4) По пункту 3 свободными остаются ровно \(t-|A|\) вершин — по одной в
компонентах \(D_i\), не сопоставленных с \(A\).
\end{proof}
\begin{graphillustration}{Пример разбиения Галлаи–Эдмондса и одно максимальное паросочетание (толстые рёбра). $D$: треугольник и две одиночные вершины ($t=3$), $A=\{a\}$, $C=\{c_1,c_2\}$. $\dfc=t-|A|=2$: свободны $d_5$ и $d_3$. Ребро $ac_1$ (между $A$ и $C$) не входит ни в одно максимальное паросочетание.}{fig:06-structure}
\begin{tikzpicture}[font=\small]
\node[gi vertex] (a) at (0,0) {$a$};
\node[gi vertex] (d1) at (-2.4,1.3) {$d_1$};
\node[gi vertex] (d2) at (-3.4,2.0) {$d_2$};
\node[gi free] (d3) at (-3.4,0.6) {$d_3$};
\node[gi vertex] (d4) at (-2.4,-0.4) {$d_4$};
\node[gi free] (d5) at (-2.4,-1.5) {$d_5$};
\node[gi vertex] (c1) at (2.2,0) {$c_1$};
\node[gi vertex] (c2) at (3.6,0) {$c_2$};
\draw[gi edge] (d1) -- (d3); \draw[gi edge] (a) -- (d1);
\draw[gi match] (d1) -- (d2);
\draw[gi edge] (d2) -- (d3);
\draw[gi match] (a) -- (d4); \draw[gi edge] (a) -- (d5);
\draw[gi edge] (a) -- (c1); \draw[gi match] (c1) -- (c2);
\node[gi group,fit=(d1)(d2)(d3)(d4)(d5),label=above:{$D(G)$}] {};
\node[gi group,fit=(a),label=above:{$A(G)$}] {};
\node[gi group,fit=(c1)(c2),label=above:{$C(G)$}] {};
\end{tikzpicture}
\end{graphillustration}

\begin{cor}
Ребро с концом в \(D(G)\) входит в некоторое максимальное паросочетание;
ребро внутри \(A\) или между \(A\) и \(C\) — ни в одно.
\end{cor}
\begin{tcolorbox}[colback=cbGrayBack,colframe=cbGray!55]
\textbf{Связь с формулой Бержа.} \(S=A(G)\) — «наилучшее» множество в формуле
Бержа: \(\odd(G-A)\ge t\) (все \(D_i\) нечётны), а \(\odd(G-A)-|A|\le\dfc(G)=t-|A|\).
\textbf{Тривиальные случаи:} есть совершенное паросочетание \(\Rightarrow D=A=\varnothing\);
\(G\) фактор-критический \(\Rightarrow D=V\). Разбиение \(D,A,C\) находится
полиномиально (алгоритм Эдмондса).
\end{tcolorbox}

\ticket{Числа Рамсея \texorpdfstring{\(R(m,n)\)}{R(m,n)}}
\label{ticket:7}
\begin{theorem}[Задача о шести знакомых]
Среди любых шести человек найдутся трое попарно знакомых или трое попарно
незнакомых. На языке графов: любой граф \(G\) на 6 вершинах содержит \(K_3\)
или независимое множество размера 3 (т.\,е. \(K_3\subseteq\overline G\)).
\end{theorem}
\begin{proof}
Возьмём вершину \(x\): \(d_G(x)+d_{\overline G}(x)=5\), поэтому одна из степеней
не меньше 3. Переходя при необходимости к \(\overline G\) (утверждение симметрично),
считаем \(d_G(x)\ge3\) и \(a,b,c\in N(x)\). Если какие-то две из \(a,b,c\) смежны,
вместе с \(x\) они образуют треугольник. Иначе \(\{a,b,c\}\) независимо.
\end{proof}
\begin{graphillustration}{Доказательство для шести вершин: у $x$ три соседа $a,b,c$. Любое ребро между ними (пунктир) замыкает треугольник с $x$; если рёбер нет — $\{a,b,c\}$ независимо.}{fig:07-six}
\begin{tikzpicture}[font=\small]
\node[gi selected] (x) at (0,1.5) {$x$};
\node[gi vertex] (a) at (-1.5,0) {$a$};
\node[gi vertex] (b) at (0,0) {$b$};
\node[gi vertex] (c) at (1.5,0) {$c$};
\node[gi vertex] (d) at (2.8,1.2) {};
\node[gi vertex] (e) at (3.3,0.2) {};
\draw[gi c1] (x) -- (a); \draw[gi c1] (x) -- (b); \draw[gi c1] (x) -- (c);
\draw[gi c2] (x) -- (d); \draw[gi c2] (x) -- (e);
\draw[gi edge,densely dotted] (a) -- (b);
\draw[gi edge,densely dotted] (b) -- (c);
\draw[gi edge,densely dotted] (a) to[bend right=30] (c);
\node[gi note,anchor=west] at (4.2,0.8) {$d_G(x)\ge3$ (сплошные);\\штриховые — рёбра $\overline G$};
\end{tikzpicture}
\end{graphillustration}


\begin{definition}[Число Рамсея]
\(R(m,n)\) — наименьшее \(k\in\N\) такое, что любой граф на \(k\) вершинах содержит
клику \(K_m\) или независимое множество размера \(n\) (\(\overline{K_n}\)).
Эквивалентно: при любой раскраске рёбер \(K_k\) в два цвета есть клика \(K_m\)
цвета 1 или клика \(K_n\) цвета 2.
\end{definition}
\begin{lemma}[Смысл определения]
Пусть \(R(m,n)=k\). Тогда (1) любой граф на \(\ge k\) вершинах содержит \(K_m\) или
независимое множество размера \(n\); (2) для любого \(t<k\) существует граф на \(t\)
вершинах без \(K_m\) и без независимого множества размера \(n\).
\end{lemma}
\begin{proof}
(1) Если \(|V|\ge k\), выберем \(V'\subseteq V\), \(|V'|=k\), и применим определение к
\(G[V']\); найденная клика или независимое множество годятся и в \(G\).
(2) При \(t=k-1\) пример существует по минимальности \(k\); при меньших \(t\) возьмём
его индуцированный подграф на \(t\) вершинах.
\end{proof}

\begin{tcolorbox}[colback=cbGrayBack,colframe=cbGray!55]
\textbf{Простейшие значения и свойства.}
\begin{itemize}\tightlist
\item \(R(m,n)=R(n,m)\): переход к дополнению меняет клики и независимые множества местами.
\item \(R(1,n)=1\): одна вершина — уже \(K_1\).
\item \(R(2,n)=n\): граф на \(n\) вершинах либо имеет ребро (\(K_2\)), либо пуст
  (независим целиком); пустой граф на \(n-1\) вершинах — контрпример.
\item Монотонность: \(R(m,n)\le R(m',n')\) при \(m\le m'\), \(n\le n'\).
\end{itemize}
\end{tcolorbox}

\begin{cor}
\(R(3,3)=6\).
\end{cor}
\begin{proof}
\(R(3,3)\le6\) по теореме о шести знакомых. Цикл \(C_5\) не содержит треугольника,
а его дополнение — снова \(C_5\) (пентаграмма), т.\,е. нет и независимого множества
размера 3. Поэтому \(R(3,3)>5\).
\end{proof}
\begin{graphillustration}{Раскраска $K_5$ без одноцветного треугольника: цвет 1 — цикл $C_5$, цвет 2 — пентаграмма (тоже $C_5$). Поэтому $R(3,3)>5$.}{fig:07-c5}
\begin{tikzpicture}[font=\small]
\foreach \i in {0,...,4} \node[gi vertex] (v\i) at (90+72*\i:1.5) {$\i$};
\foreach \i [evaluate=\i as \j using {int(mod(\i+1,5))}] in {0,...,4} \draw[gi c1] (v\i) -- (v\j);
\foreach \i [evaluate=\i as \j using {int(mod(\i+2,5))}] in {0,...,4} \draw[gi c2] (v\i) -- (v\j);
\node[gi note,anchor=west] at (2.3,0.4) {сплошные: $G=C_5$\\штриховые: $\overline G\cong C_5$};
\end{tikzpicture}
\end{graphillustration}

Конечность \(R(m,n)\) для всех \(m,n\) будет доказана в билете~\ref{ticket:8}.
Известно немного точных значений: \(R(3,4)=9\), \(R(3,5)=14\), \(R(4,4)=18\);
уже \(R(5,5)\) неизвестно (\(43\le R(5,5)\le46\)).

\ticket{Теорема Эрдёша–Секереша}
\label{ticket:8}
\begin{theorem}[Эрдёш–Секереш, 1935]
При \(m,n\ge2\)
\[
R(m,n)\le R(m-1,n)+R(m,n-1).
\]
В частности, все числа \(R(m,n)\) конечны (индукция по \(m+n\) с базой \(R(1,n)=R(m,1)=1\)).
\end{theorem}
\begin{proof}
Пусть \(G\) — граф на \(N=R(m-1,n)+R(m,n-1)\) вершинах, \(v\in V\). Положим
\[
A=N(v),\qquad B=V\setminus(N(v)\cup\{v\}),\qquad |A|+|B|=N-1 .
\]
Значит, \(|A|\ge R(m-1,n)\) или \(|B|\ge R(m,n-1)\) (иначе \(|A|+|B|\le N-2\)).
\begin{itemize}\tightlist
\item \(|A|\ge R(m-1,n)\): в \(G[A]\) есть \(K_{m-1}\) — добавив \(v\) (смежна со всеми
  вершинами \(A\)), получаем \(K_m\); либо есть независимое множество размера \(n\).
\item \(|B|\ge R(m,n-1)\): в \(G[B]\) есть \(K_m\); либо независимое множество размера
  \(n-1\), и вместе с \(v\) (не смежна ни с одной вершиной \(B\)) — размера \(n\).
\end{itemize}
По лемме из билета~\ref{ticket:7} этого достаточно.
\end{proof}
\begin{graphillustration}{Разбиение в доказательстве Эрдёша–Секереша: $v$ смежна со всеми вершинами $A$ и ни с одной вершиной $B$. Клика $K_{m-1}$ в $A$ дополняется вершиной $v$ до $K_m$; независимое множество размера $n-1$ в $B$ — до размера $n$.}{fig:08-split}
\begin{tikzpicture}[font=\small]
\node[gi selected] (v) at (0,0) {$v$};
\node[draw,ellipse,minimum width=2.8cm,minimum height=1.8cm] (A) at (-2.8,-1.8) {};
\node[draw,ellipse,minimum width=2.8cm,minimum height=1.8cm] (B) at (2.8,-1.8) {};
\node[gi note] at (-2.8,-3.0) {$A=N(v)$, $|A|\ge R(m-1,n)$?};
\node[gi note] at (2.8,-3.0) {$B=V\setminus N[v]$, $|B|\ge R(m,n-1)$?};
\foreach \x/\y [count=\i] in {-3.6/-1.5,-2.8/-1.2,-2.0/-1.5,-3.2/-2.2,-2.4/-2.2}
  {\node[gi small] (a\i) at (\x,\y) {}; \draw[gi c1] (v) -- (a\i);}
\foreach \x/\y [count=\i] in {2.0/-1.5,2.8/-1.2,3.6/-1.5,2.4/-2.2,3.2/-2.2}
  {\node[gi small] (b\i) at (\x,\y) {}; \draw[gi c2,black!45] (v) -- (b\i);}
\node[gi note,anchor=west] at (5.0,0) {сплошные — рёбра $G$\\штриховые — рёбра $\overline G$};
\end{tikzpicture}
\end{graphillustration}


\begin{theorem}[Усиление при чётных слагаемых]
Если \(R(m-1,n)\) и \(R(m,n-1)\) оба чётны, то \(R(m,n)\le R(m-1,n)+R(m,n-1)-1\).
\end{theorem}
\begin{proof}
Пусть \(N=R(m-1,n)+R(m,n-1)-1\) (нечётно) и в \(G\) на \(N\) вершинах нет ни \(K_m\),
ни независимого множества размера \(n\). Как и выше, для каждой вершины
\(|A|\le R(m-1,n)-1\) и \(|B|\le R(m,n-1)-1\), а \(|A|+|B|=N-1\) — значит, обе оценки
точные и \(d(v)=|A|=R(m-1,n)-1\) нечётно для всех \(v\). Нечётное число вершин
нечётной степени противоречит лемме о рукопожатиях.
\end{proof}
Пример: \(R(3,3)\le R(2,3)+R(3,2)=6\); \(R(3,4)\le R(2,4)+R(3,3)-1=4+6-1=9\)
(обе величины чётны) — точное значение.

\begin{cor}[Биномиальная оценка]
\[
R(m,n)\le\binom{m+n-2}{n-1}=\binom{m+n-2}{m-1}.
\]
\end{cor}
\begin{proof}
Индукция по \(m+n\). Если \(m=1\) или \(n=1\): \(R=1\le\binom{m+n-2}{n-1}\)
(при \(m=1\) биномиальный коэффициент \(\binom{n-1}{n-1}=1\), при \(n=1\) — \(\binom{m-1}{0}=1\)).
При \(m,n\ge2\) по теореме и тождеству Паскаля
\[
R(m,n)\le\binom{m+n-3}{n-1}+\binom{m+n-3}{n-2}=\binom{m+n-2}{n-1}.\qedhere
\]
\end{proof}
\begin{cor}
\(R(n,n)\le\binom{2n-2}{n-1}\sim\dfrac{4^{n-1}}{\sqrt{\pi n}}=O\!\left(\dfrac{4^n}{\sqrt n}\right)\).
\end{cor}
\begin{proof}
Формула Стирлинга: \(\binom{2k}{k}\sim\frac{4^k}{\sqrt{\pi k}}\) при \(k=n-1\).
\end{proof}

\ticket{Нижняя оценка Эрдёша и её усиления}
\label{ticket:9}
\begin{lemma}[Простая нижняя оценка (лекции)]
При \(3\le m\le n\): \(R(m,n)\ge m+n-1\).
\end{lemma}
\begin{proof}
Граф \(K_{m-1,n-1}\) имеет \(m+n-2\) вершин. Клик размера 3 нет (две из трёх вершин
попали бы в одну долю), а \(m\ge3\) — значит, нет и \(K_m\). Независимое множество лежит
в одной доле (вершины разных долей смежны), поэтому \(\alpha=n-1<n\).
\end{proof}
Конструкция «\(m-1\) непересекающихся клик \(K_{n-1}\)» даёт заметно лучшее
\(R(m,n)>(m-1)(n-1)\) (см. рис.~\ref{fig:12-chvatal} в билете~\ref{ticket:12}), но всё
ещё лишь полиномиальную оценку. Эрдёш показал, что рост экспоненциальный.
\begin{graphillustration}{Простая нижняя оценка при $m=3$, $n=4$: в $K_{2,3}$ (5 вершин) нет треугольника, а $\alpha=3<4$. Значит, $R(3,4)\ge6$ (истинное значение 9).}{fig:09-bipartite}
\begin{tikzpicture}[font=\small]
\node[gi vertex] (a1) at (0.8,1.2) {};
\node[gi vertex] (a2) at (2.2,1.2) {};
\node[gi vertex] (b1) at (0,0) {};
\node[gi vertex] (b2) at (1.5,0) {};
\node[gi vertex] (b3) at (3,0) {};
\foreach \a in {a1,a2} \foreach \b in {b1,b2,b3} \draw[gi edge] (\a) -- (\b);
\node[gi group,fit=(a1)(a2),label=right:{$m-1=2$}] {};
\node[gi group,fit=(b1)(b2)(b3),label=right:{$n-1=3$}] {};
\end{tikzpicture}
\end{graphillustration}


\begin{theorem}[Эрдёш, 1947]
При \(n\ge2\): \(R(n,n)\ge2^{n/2}\).
\end{theorem}
\begin{proof}
\(R(2,2)=2=2^{2/2}\); пусть \(n\ge3\). Зафиксируем \(\nu<2^{n/2}\) помеченных вершин и
покажем, что есть граф на них без \(K_n\) и без независимого множества размера \(n\).
Пусть \(\mathcal G(\nu)=2^{\binom\nu2}\) — число всех графов на этих вершинах,
\(T(\nu,n)\) — число графов, содержащих \(K_n\), \(P(\nu,n)\) — содержащих независимое
множество размера \(n\). Если \(T+P<\mathcal G\), искомый граф существует.

\emph{\(T=P\)}: дополнение — биекция между этими классами.

\emph{Оценка \(T\)}: выберем \(n\) вершин будущей клики (\(\binom{\nu}{n}\) способов), рёбра
между ними проведём, остальные \(\binom\nu2-\binom n2\) пар — произвольно. Каждый граф с
\(K_n\) учтён хотя бы раз:
\[
T(\nu,n)\le\binom{\nu}{n}2^{\binom\nu2-\binom n2}
\le\frac{\nu^n}{n!}\,2^{-\frac{n(n-1)}2}\,\mathcal G(\nu)
<\frac{2^{n^2/2}}{n!}\,2^{-\frac{n^2}2+\frac n2}\,\mathcal G(\nu)
=\frac{(\sqrt2)^n}{n!}\,\mathcal G(\nu).
\]
\emph{\((\sqrt2)^n/n!<\tfrac12\) при \(n\ge3\)}: при \(n=3\) это \(2\sqrt2/6\approx0{,}47\), а при
переходе \(n\to n+1\) выражение умножается на \(\sqrt2/(n+1)<1\).

Итак, \(T<\mathcal G/2\), \(P=T<\mathcal G/2\), \(T+P<\mathcal G\). Значит, \(R(n,n)>\nu\) для каждого
целого \(\nu<2^{n/2}\), откуда \(R(n,n)\ge2^{n/2}\).
\end{proof}
\begin{tcolorbox}[colback=cbGrayBack,colframe=cbGray!55]
\textbf{Вероятностный смысл.} Выберем граф на \(\nu\) вершинах случайно (каждое ребро
независимо с вероятностью \(\tfrac12\)). Доказано: вероятность увидеть \(K_n\) или
независимое \(n\)-множество меньше 1, значит, хороший граф существует, хотя
явно его построить никто не умеет.
\end{tcolorbox}
\begin{graphillustration}{Подсчёт в доказательстве Эрдёша: фиксируем $n$ вершин и проводим все $\binom n2$ рёбер между ними (сплошные); остальные пары (серые) выбираем произвольно. Получаем не больше $\binom\nu n2^{\binom\nu2-\binom n2}$ графов с кликой $K_n$.}{fig:09-counting}
\begin{tikzpicture}[font=\small]
\foreach \i in {0,...,7} \node[gi small] (v\i) at (90+45*\i:1.4) {};
\foreach \i in {0,...,7} \foreach \j in {0,...,7} {\ifnum\i<\j \draw[black!18] (v\i) -- (v\j); \fi}
\draw[gi c1] (v0) -- (v1) -- (v2) -- (v0);
\draw[gi c1] (v0) -- (v7) -- (v1); \draw[gi c1] (v2) -- (v7);
\node[gi group,fit=(v0)(v1)(v2)(v7),label=above:{$n$ вершин}] {};
\node[gi note,anchor=west] at (2.2,0) {$\binom\nu n$ выборов клики,\\$2^{\binom\nu2-\binom n2}$ — остальное};
\end{tikzpicture}
\end{graphillustration}


\subsection*{Усиления}
\begin{cor}
При \(2\le m\le n\): \(R(m,n)\ge R(m,m)\ge2^{m/2}\), т.\,е. \(R(m,n)\ge2^{\min(m,n)/2}\).
\end{cor}
\begin{cor}[Общий критерий из доказательства]
Если \(\binom{\nu}{n}2^{1-\binom n2}<1\), то \(R(n,n)>\nu\).
\end{cor}
\begin{proof}
Из этого условия \(T+P\le2\binom\nu n2^{\binom\nu2-\binom n2}<\mathcal G(\nu)\).
\end{proof}
\begin{cor}[Точная форма оценки Эрдёша]
\[
R(n,n)\ge\bigl(1+o(1)\bigr)\frac{n}{e\sqrt2}\,2^{n/2}.
\]
\end{cor}
\begin{proof}
Используем \(\binom\nu n\le\frac{\nu^n}{n!}\le\left(\frac{e\nu}{n}\right)^n\) (так как \(n!\ge(n/e)^n\)):
\[
\binom{\nu}{n}2^{1-\binom n2}\le2\left(\frac{e\nu}{n}\right)^n2^{-\frac{n^2}{2}+\frac n2}
=2\left(\frac{e\sqrt2\,\nu}{n\,2^{n/2}}\right)^n .
\]
Для фиксированного \(\varepsilon>0\) и \(\nu=\bigl\lfloor(1-\varepsilon)\frac{n}{e\sqrt2}2^{n/2}\bigr\rfloor\)
правая часть не больше \(2(1-\varepsilon)^n\to0\), т.\,е. при больших \(n\) меньше 1, и
\(R(n,n)>\nu\) по предыдущему следствию.
\end{proof}
\textbf{Итог.} \(\sqrt2^{\,n}\lesssim R(n,n)\lesssim4^n\); показатели \(\sqrt2\) и \(4\) с 1947 года
улучшены лишь незначительно (верхняя — в 2023 году до \((4-\varepsilon)^n\)).

\ticket{Несколько цветов. Теорема Гринвуда–Глисона}
\label{ticket:10}
\begin{definition}[Многоцветные числа Рамсея]
\(R(n_1,\ldots,n_k)\) — наименьшее \(t\) такое, что при любой раскраске рёбер \(K_t\)
в \(k\) цветов \(\varphi\colon E(K_t)\to\{1,\ldots,k\}\) найдётся \(i\) и клика \(K_{n_i}\),
все рёбра которой цвета \(i\).
\end{definition}
При \(k=2\) получаем \(R(n_1,n_2)\): граф \((V,\varphi^{-1}(1))\) содержит \(K_{n_1}\),
а граф \((V,\varphi^{-1}(2))\) — его дополнение — \(K_{n_2}\).
\begin{tcolorbox}[colback=cbGrayBack,colframe=cbGray!55]
\textbf{Свойства.} \(R\) симметрично (не меняется при перестановке аргументов) и
монотонно по каждому аргументу. Если некоторое \(n_i=1\), то \(R=1\) (одна вершина —
клика \(K_1\) любого цвета). Если \(n_i=2\), то цвет \(i\) нельзя использовать вовсе, и
\(R(n_1,\ldots,2,\ldots,n_k)=R(n_1,\ldots,\widehat{n_i},\ldots,n_k)\).
\end{tcolorbox}

\begin{theorem}[Гринвуд–Глисон, 1955]
При \(n_1,\ldots,n_k\ge2\)
\[
R(n_1,\ldots,n_k)\le\sum_{i=1}^{k}R(n_1,\ldots,n_i-1,\ldots,n_k)-k+2 .
\]
В лекциях записано более слабое \(R(n_1,\ldots,n_k)\le\sum_iR(n_1,\ldots,n_i-1,\ldots,n_k)\).
\end{theorem}
\begin{proof}
Обозначим \(R_i=R(n_1,\ldots,n_i-1,\ldots,n_k)\) и рассмотрим раскраску \(K_N\),
\(N=\sum R_i-k+2\). Возьмём вершину \(v\); пусть \(A_i\) — множество вершин, соединённых
с \(v\) рёбрами цвета \(i\). Тогда \(\sum|A_i|=N-1=\sum(R_i-1)+1\), и по принципу
Дирихле \(|A_i|\ge R_i\) для некоторого \(i\). По определению \(R_i\) в клике на \(A_i\) есть
\begin{itemize}\tightlist
\item либо клика \(K_{n_i-1}\) цвета \(i\) — с вершиной \(v\) она даёт \(K_{n_i}\) цвета \(i\);
\item либо для некоторого \(j\ne i\) клика \(K_{n_j}\) цвета \(j\).
\end{itemize}
В обоих случаях нужная одноцветная клика найдена.
\end{proof}
\begin{graphillustration}{Шаг доказательства Гринвуда–Глисона при $k=3$: рёбра из $v$ делятся по цветам на $A_1,A_2,A_3$. Хотя бы одно $|A_i|\ge R_i$; внутри него ищем $K_{n_i-1}$ цвета $i$ (с $v$ — $K_{n_i}$) или $K_{n_j}$ другого цвета.}{fig:10-pigeon}
\begin{tikzpicture}[font=\small]
\node[gi selected] (v) at (0,0) {$v$};
\foreach \k/\ang/\sty/\lab in {1/150/gi c1/A_1,2/30/gi c2/A_2,3/270/gi c3/A_3} {
  \node[draw,ellipse,minimum width=2.2cm,minimum height=1.2cm] (E\k) at (\ang:2.6) {};
  \node[gi note] at ($(\ang:2.6)+(1.45,0.45)$) {$\lab$};
  \foreach \d [count=\j] in {-0.6,0,0.6} {\node[gi small] (p\k-\j) at ($(\ang:2.6)+(\d,-0.1)$) {}; \draw[\sty] (v) -- (p\k-\j);}
}
\node[gi note,anchor=west] at (4.2,-1.4) {$\sum|A_i|=N-1=\sum(R_i-1)+1$\\$\Rightarrow\ \exists i:\ |A_i|\ge R_i$};
\end{tikzpicture}
\end{graphillustration}

\begin{cor}
Все числа \(R(n_1,\ldots,n_k)\) конечны.
\end{cor}

\begin{cor}[Мультиномиальная оценка]
\[
R(n_1,\ldots,n_k)\le\binom{n_1+\cdots+n_k-k}{n_1-1,\ \ldots,\ n_k-1}
=\frac{(n_1+\cdots+n_k-k)!}{(n_1-1)!\cdots(n_k-1)!}.
\]
\end{cor}
\begin{proof}
Индукция по \(n_1+\cdots+n_k\). Если какое-то \(n_i=1\), левая часть равна 1, а
мультиномиальный коэффициент не меньше 1. Иначе все \(a_i=n_i-1\ge1\), и
по теореме (в слабой форме) и тождеству для мультиномиальных коэффициентов
\[
\binom{s}{a_1,\ldots,a_k}=\sum_{i=1}^{k}\binom{s-1}{a_1,\ldots,a_i-1,\ldots,a_k},
\qquad s=a_1+\cdots+a_k,
\]
получаем \(R(n_1,\ldots,n_k)\le\sum_iR(\ldots,n_i-1,\ldots)\le\binom{s}{a_1,\ldots,a_k}\).
Тождество — разбор по тому, какой «тип» стоит на первом месте в слове из
\(a_1\) букв первого типа, \ldots, \(a_k\) букв \(k\)-го типа.
\end{proof}
При \(k=2\) это ровно \(\binom{m+n-2}{m-1}\) из билета~\ref{ticket:8}.
\begin{example}
\(R(3,3,3)\le R(2,3,3)+R(3,2,3)+R(3,3,2)-3+2=3\cdot6-1=17\) (здесь
\(R(2,3,3)=R(3,3)=6\)). Это точное значение: \(R(3,3,3)=17\).
Мультиномиальная оценка даёт \(\frac{6!}{2!\,2!\,2!}=90\).
\end{example}

\subsection*{Применение: теорема Шура}
\begin{theorem}[Шур, 1916]
Для любого \(k\ge1\) существует \(n\) такое, что при любой раскраске чисел
\(1,\ldots,n\) в \(k\) цветов найдутся \(x,y,z\) одного цвета с \(x+y=z\).
\end{theorem}
\begin{proof}
Пусть \(N=R(\underbrace{3,\ldots,3}_k)\), \(n=N-1\), и \(\varphi\colon\{1,\ldots,n\}\to\{1,\ldots,k\}\) —
раскраска чисел. Раскрасим рёбра \(K_N\) на вершинах \(1,\ldots,N\) правилом
\(\psi(uv)=\varphi(|u-v|)\). По определению \(N\) есть одноцветный треугольник
\(a<b<c\): \(\varphi(b-a)=\varphi(c-b)=\varphi(c-a)\). Положим \(x=b-a\), \(y=c-b\),
\(z=c-a\): тогда \(x+y=z\), и все три одного цвета (допускается \(x=y\)).
\end{proof}
\begin{graphillustration}{Сведение Шура к числам Рамсея: ребро $uv$ получает цвет числа $|u-v|$. Одноцветный треугольник $a<b<c$ даёт $x=b-a$, $y=c-b$, $z=c-a$ одного цвета, $x+y=z$. Справа — раскраска $1,\ldots,4$ в два цвета без решения (синий контур — цвет 1, двойной серый — цвет 2): оценка $n=R(3,3)-1=5$ точна.}{fig:10-schur}
\begin{tikzpicture}[font=\small]
\node[gi vertex] (a) at (0,0) {$a$};
\node[gi vertex] (b) at (2.4,0) {$b$};
\node[gi vertex] (c) at (1.2,1.8) {$c$};
\draw[gi c1] (a) -- node[below] {$b-a$} (b);
\draw[gi c1] (b) -- node[right] {$c-b$} (c);
\draw[gi c1] (a) -- node[left] {$c-a$} (c);
\begin{scope}[shift={(5.0,0.8)}]
\node[gi vertex,draw=cbBlue,line width=1.4pt] at (0,0) {$1$};
\node[gi selected] at (0.9,0) {$2$};
\node[gi selected] at (1.8,0) {$3$};
\node[gi vertex,draw=cbBlue,line width=1.4pt] at (2.7,0) {$4$};
\node[gi note,anchor=west] at (3.3,0) {$1+1=2$, $1+3=4$, $2+2=4$\\— везде разные цвета};
\end{scope}
\end{tikzpicture}
\end{graphillustration}


\ticket{Графовые числа Рамсея}
\label{ticket:11}
\begin{definition}[Графовое число Рамсея]
Для графов \(G_1,G_2\) число \(R(G_1,G_2)\) — наименьшее \(t\) такое, что при любой
раскраске рёбер \(K_t\) в два цвета найдётся подграф цвета 1, изоморфный \(G_1\),
или подграф цвета 2, изоморфный \(G_2\). Подграф \emph{не обязан быть индуцированным}:
нужны лишь рёбра \(G_i\) нужного цвета, остальные рёбра на тех же вершинах любые.
Аналогично для \(k\) цветов: \(R(G_1,\ldots,G_k)\).
\end{definition}
\begin{theorem}[Существование]
\[
R(G_1,\ldots,G_k)\le R\bigl(|G_1|,\ldots,|G_k|\bigr),
\]
где \(|G_i|=v(G_i)\). В частности, графовые числа Рамсея конечны.
\end{theorem}
\begin{proof}
Пусть \(t=R(|G_1|,\ldots,|G_k|)\) и дана раскраска \(K_t\). По определению найдётся \(i\) и
клика \(K_{|G_i|}\) цвета \(i\). Граф \(G_i\) на \(|G_i|\) вершинах — подграф полного графа
на стольких же вершинах, поэтому в этой клике есть копия \(G_i\), все рёбра которой
цвета \(i\). Конечность следует из билета~\ref{ticket:10}.
\end{proof}
\begin{cor}
\(R(K_{n_1},\ldots,K_{n_k})=R(n_1,\ldots,n_k)\): для клик неравенство превращается в
определение.
\end{cor}
\begin{graphillustration}{Слева: одноцветная клика $K_4$ содержит любой граф на 4 вершинах, например цикл $C_4$ (толстые рёбра); поэтому $R(C_4,\cdot)\le R(4,\cdot)$. Справа: в любой 2-раскраске $K_3$ два ребра одного цвета образуют $P_3$.}{fig:11-embed}
\begin{tikzpicture}[font=\small]
\begin{scope}
\node[gi vertex] (a) at (0,0) {};
\node[gi vertex] (b) at (1.6,0) {};
\node[gi vertex] (c) at (1.6,1.6) {};
\node[gi vertex] (d) at (0,1.6) {};
\draw[gi c1] (a) -- (c); \draw[gi c1] (b) -- (d);
\draw[gi match] (a) -- (b) -- (c) -- (d) -- (a);
\node[gi note] at (0.8,-0.6) {$C_4\subseteq K_4$ цвета 1};
\end{scope}
\begin{scope}[shift={(4.2,0)}]
\node[gi vertex] (x) at (0,0) {};
\node[gi vertex] (y) at (1.8,0) {};
\node[gi vertex] (z) at (0.9,1.5) {};
\draw[gi c1,line width=2pt] (x) -- (z); \draw[gi c1,line width=2pt] (z) -- (y);
\draw[gi c2] (x) -- (y);
\node[gi note] at (0.9,-0.6) {одноцветный $P_3$};
\end{scope}
\end{tikzpicture}
\end{graphillustration}

\begin{example}[Графовые числа бывают много меньше]
Несколько простых значений:
\begin{itemize}\tightlist
\item \(R(P_3,P_3)=3\) (\(P_3\) — путь из двух рёбер), тогда как \(R(3,3)=6\).
  В \(K_3\) из трёх рёбер два одного цвета, а любые два ребра треугольника
  имеют общую вершину — одноцветный \(P_3\). В \(K_2\) одно ребро, \(P_3\) нет.
\item \(R(K_2,G)=v(G)\) при \(v(G)\ge2\): раскраска без ребра цвета 1 — сплошь цвета 2, и
  \(G\subseteq K_t\iff t\ge v(G)\).
\item \(R(T_n,K_m)=(m-1)(n-1)+1\) для любого дерева \(T_n\) — билет~\ref{ticket:12}.
\end{itemize}
\end{example}
\begin{tcolorbox}[colback=cbOrangeBack,colframe=cbOrange!45!cbGray,
  borderline west={1.6pt}{0pt}{cbOrange}]
\textbf{Индуцированный случай (только формулировка).} Граф \(G\) называется
\(H\)-рамсеевским, если при любой раскраске рёбер \(G\) в два цвета есть одноцветный
\emph{индуцированный} подграф \(G\), изоморфный \(H\). Для любого \(H\) такой \(G\)
существует (индуцированная теорема Рамсея: Деуб, Эрдёш–Хайнал–Поша, Рёдль, 1970-е),
но доказательство значительно сложнее и в программу не входит. Для клик
«индуцированный» ничего не меняет — это обычная теорема Рамсея.
\end{tcolorbox}

\ticket{Число Рамсея для дерева и клики}
\label{ticket:12}
\begin{theorem}[Хватал, 1977]
Для любого дерева \(T_n\) на \(n\) вершинах и \(m\ge1\)
\[
R(T_n,K_m)=(m-1)(n-1)+1 .
\]
Здесь ищется дерево \(T_n\) цвета 1 или клика \(K_m\) цвета 2. Ответ не зависит от формы дерева!
\end{theorem}
\textit{Нижняя оценка: \(R(T_n,K_m)\ge(m-1)(n-1)+1\).}
Раскрасим рёбра \(K_{(m-1)(n-1)}\): разобьём вершины на \(m-1\) групп по \(n-1\) вершине,
рёбра внутри групп — цвета 1, между группами — цвета 2. Компоненты связности
графа цвета 1 имеют по \(n-1\) вершине, поэтому связного \(T_n\) цвета 1 нет. Граф
цвета 2 — полный \((m-1)\)-дольный: среди любых \(m\) вершин две в одной доле, они не
соединены ребром цвета 2, так что \(K_m\) цвета 2 нет.
\begin{graphillustration}{Экстремальная раскраска при $m=3$, $n=4$: две группы по 3 вершины. Внутри групп цвет 1 (сплошные) — компоненты из $n-1=3$ вершин, $T_4$ не помещается; между группами цвет 2 (штриховые) — двудольный граф без треугольника.}{fig:12-chvatal}
\begin{tikzpicture}[font=\small]
\node[gi vertex] (a1) at (0,1.4) {};
\node[gi vertex] (a2) at (-0.8,0) {};
\node[gi vertex] (a3) at (0.8,0) {};
\node[gi vertex] (b1) at (4,1.4) {};
\node[gi vertex] (b2) at (3.2,0) {};
\node[gi vertex] (b3) at (4.8,0) {};
\foreach \x in {a1,a2,a3} \foreach \y in {b1,b2,b3} \draw[gi c2] (\x) -- (\y);
\draw[gi c1] (a1) -- (a2) -- (a3) -- (a1);
\draw[gi c1] (b1) -- (b2) -- (b3) -- (b1);
\node[gi group,fit=(a1)(a2)(a3),label=below:{$K_{n-1}$}] {};
\node[gi group,fit=(b1)(b2)(b3),label=below:{$K_{n-1}$}] {};
\node[gi note,anchor=west] at (5.6,0.7) {$(m-1)(n-1)=6$ вершин:\\нет $T_4$ цвета 1,\\нет $K_3$ цвета 2};
\end{tikzpicture}
\end{graphillustration}


\begin{lemma}\label{lem:trees}
Пусть \(m>1\), \(H\) — граф с \(v(H)\ge(m-1)(n-1)+1\) и \(\alpha(H)\le m-1\). Тогда \(H\)
содержит (как подграф) любое дерево на \(n\) вершинах.
\end{lemma}
\begin{proof}
Индукция по \(n\) при фиксированном \(m\). База \(n=1\): в \(H\) есть хотя бы одна вершина.

Переход \(n-1\to n\). Пусть \(u\) — лист \(T_n\), \(a\) — его сосед, \(T_{n-1}=T_n-u\).
Возьмём максимальное независимое множество \(U\) графа \(H\): \(|U|=\alpha(H)\le m-1\). Тогда
\[
v(H-U)\ge(m-1)(n-1)+1-(m-1)=(m-1)(n-2)+1,\qquad \alpha(H-U)\le\alpha(H)\le m-1,
\]
и по предположению индукции \(H-U\) содержит копию \(T_{n-1}\). Пусть \(a'\) — образ \(a\).
Множество \(U\cup\{a'\}\) не независимо (иначе \(U\) не максимально), поэтому \(a'\) смежна
с некоторой \(x\in U\). Вершина \(x\) не входит в копию \(T_{n-1}\subseteq H-U\);
присоединив её к \(a'\), получаем копию \(T_n\).
\end{proof}
\begin{graphillustration}{Шаг индукции в лемме: копия $T_{n-1}$ найдена в $H-U$; вершина $a'$ (образ соседа удалённого листа) смежна с некоторой $x\in U$ по максимальности $U$. Ребро $a'x$ (толстое) достраивает $T_n$.}{fig:12-step}
\begin{tikzpicture}[font=\small]
\node[draw,rounded corners,minimum width=5.2cm,minimum height=2.4cm] (HU) at (0,0) {};
\node[gi note,anchor=north west] at (-2.55,1.15) {$H-U$};
\node[gi vertex] (t1) at (-1.6,-0.2) {};
\node[gi vertex] (t2) at (-0.6,0.4) {};
\node[gi vertex] (t3) at (-0.6,-0.7) {};
\node[gi vertex] (t4) at (0.5,0.4) {};
\node[gi selected] (a) at (1.6,0.4) {$a'$};
\draw[gi edge] (t1) -- (t2) -- (t4) -- (a); \draw[gi edge] (t1) -- (t3);
\node[draw,rounded corners,minimum width=1.4cm,minimum height=2.4cm] (U) at (4.2,0) {};
\node[gi note,anchor=north] at (4.2,1.6) {$U$, $|U|\le m-1$};
\node[gi vertex] (x) at (4.2,0.6) {$x$};
\node[gi vertex] (y) at (4.2,-0.6) {};
\draw[gi match] (a) -- (x);
\end{tikzpicture}
\end{graphillustration}


\textit{Верхняя оценка.} Рассмотрим раскраску \(K_N\), \(N=(m-1)(n-1)+1\), и граф \(H\)
рёбер цвета 1. Если нет \(K_m\) цвета 2, то в \(H\) нет независимого множества размера
\(m\): \(\alpha(H)\le m-1\) (при \(m=1\) утверждение тривиально: \(N=1\) и \(K_1\) есть). По
лемме~\ref{lem:trees} \(H\) содержит \(T_n\), т.\,е. есть \(T_n\) цвета 1. \qed

\begin{example}
\(R(P_3,K_3)=2\cdot2+1=5\), \(R(K_{1,3},K_3)=R(P_4,K_3)=7\): звезда и путь на 4 вершинах
дают одно и то же число.
\end{example}

\ticket{Интервальные графы и графы сравнимости}
\label{ticket:13}
\subsection*{Совершенные графы}
\begin{definition}[Совершенный граф]
Граф \(G\) \emph{совершенный}, если \(\chi(H)=\omega(H)\) для каждого индуцированного
подграфа \(H\le G\). Индуцированный подграф совершенного графа совершенен.
\end{definition}
Всегда \(\chi\ge\omega\). Нечётный цикл \(C_{2k+1}\), \(k\ge2\), не совершенен:
\(\chi=3\), \(\omega=2\). Совершенны: \(K_n\) (\(\chi=\omega=n\)), пустые графы (\(\chi=\omega=1\)),
двудольные графы (\(\chi=\omega=2\) при наличии рёбер; индуцированные подграфы тоже
двудольны), в частности чётные циклы.

\subsection*{Интервальные графы}
\begin{definition}[Интервальный граф]
\(G\) — \emph{интервальный}, если существуют интервалы \(I_1,\ldots,I_n\subset\R\) такие, что
\(v_iv_j\in E\iff I_i\cap I_j\ne\varnothing\). Далее \(I_k=(\ell_k,r_k)\).
\end{definition}
\begin{graphillustration}{Пример интервального графа. Пересекающиеся интервалы дают рёбра; вертикальная прямая пересекает 3 интервала — это $\omega=\chi=3$ (клика $\{2,3,4\}$). Цифра в вершине — номер интервала.}{fig:13-intervals}
\begin{tikzpicture}[font=\small]
\draw[gi interval] (0,2.0) -- (2.0,2.0) node[right] {$I_1$};
\draw[gi interval] (1.2,1.4) -- (3.6,1.4) node[right] {$I_2$};
\draw[gi interval] (2.6,0.8) -- (5.0,0.8) node[right] {$I_3$};
\draw[gi interval] (3.0,0.2) -- (4.2,0.2) node[right] {$I_4$};
\draw[gi interval] (4.6,1.6) -- (6.2,1.6) node[right] {$I_5$};
\draw[->] (-0.3,-0.4) -- (6.9,-0.4) node[right] {$\R$};
\draw[gi removed] (3.3,-0.6) -- (3.3,2.3);
\begin{scope}[shift={(9,1)}]
\node[gi vertex] (1) at (-1.6,0.6) {1};
\node[gi vertex] (2) at (-0.6,0.6) {2};
\node[gi vertex] (3) at (0.6,0.6) {3};
\node[gi vertex] (4) at (0,-0.5) {4};
\node[gi vertex] (5) at (1.6,0.6) {5};
\draw[gi edge] (1) -- (2) -- (3) -- (5); \draw[gi edge] (2) -- (4) -- (3);
\end{scope}
\end{tikzpicture}
\end{graphillustration}

\begin{lemma}[Индуцированные подграфы]\label{lem:int-induced}
Индуцированный подграф интервального графа интервален: достаточно оставить
интервалы его вершин.
\end{lemma}
\begin{lemma}[Свойство Хелли на прямой; \(\omega\)]\label{lem:helly}
Если интервалы попарно пересекаются, то у них есть общая точка. Поэтому
\(\omega(G)\) равно наибольшему числу интервалов, покрывающих одну точку.
\end{lemma}
\begin{proof}
Для \(S\subseteq V\): \(\bigcap_{k\in S}I_k=(\max\ell_k,\min r_k)\). Пусть \(\ell_i=\max\ell_k\),
\(r_j=\min r_k\). Из \(I_i\cap I_j\ne\varnothing\) следует \(\max(\ell_i,\ell_j)<\min(r_i,r_j)\), т.\,е.
\(\ell_i<r_j\), и пересечение непусто. Значит, клика — это множество интервалов с общей
точкой, и наоборот: интервалы через одну точку попарно пересекаются.
\end{proof}
\begin{lemma}[\(\chi\) интервального графа]\label{lem:greedy-int}
\(\chi(G)\) равно наибольшему числу интервалов, покрывающих одну точку, т.\,е. \(\chi=\omega\).
\end{lemma}
\begin{proof}
Жадный алгоритм: перебираем интервалы по возрастанию левых концов и даём
каждому наименьший цвет, не занятый уже покрашенными соседями. Пусть
обрабатывается \(I=(\ell,r)\), \(J=(\ell_J,r_J)\) — ранее покрашенный сосед. Тогда
\(\ell_J\le\ell\) и \(J\cap I\ne\varnothing\) влечёт \(r_J>\ell\). Значит, все такие \(J\) вместе с \(I\)
содержат точки \((\ell,\ell+\varepsilon)\) при малом \(\varepsilon\): это клика, и соседей не больше
\(\omega-1\). Цвет \(\le\omega\), поэтому \(\chi\le\omega\le\chi\).
\end{proof}
\begin{theorem}
Интервальный граф совершенен.
\end{theorem}
\begin{proof}
Каждый индуцированный подграф интервален (лемма~\ref{lem:int-induced}), и для него \(\chi=\omega\) (лемма~\ref{lem:greedy-int}).
\end{proof}
\begin{example}[\(C_n\), \(n\ge4\), не интервален]
Более того, в интервальном графе нет индуцированного цикла длины \(\ge4\). Пусть
\(v_1v_2\ldots v_k\) (\(k\ge4\)) — индуцированный цикл и, не умаляя общности, \(r_1\) —
наименьший правый конец среди его интервалов. Соседи \(v_2,v_k\) пересекают \(I_1\) и
кончаются правее \(r_1\): \(r_2,r_k\ge r_1\), \(\ell_2,\ell_k<r_1\). Тогда точки чуть левее \(r_1\)
лежат и в \(I_2\), и в \(I_k\) — \(v_2v_k\in E\), хорда. Противоречие.
\end{example}
\begin{graphillustration}{Почему цикл длины $\ge4$ не интервален: у $I_1$ наименьший правый конец $r_1$. Оба соседа $v_2,v_k$ пересекают $I_1$ и продолжаются правее $r_1$, поэтому пересекаются между собой около $r_1$ (серая полоса) — появляется хорда $v_2v_k$.}{fig:13-c4}
\begin{tikzpicture}[font=\small]
\fill[black!12] (2.7,-0.2) rectangle (3.0,1.6);
\draw[gi interval] (0.5,1.3) -- (3.0,1.3) node[right] {$I_1$};
\draw[gi interval] (1.5,0.75) -- (4.8,0.75) node[right] {$I_2$};
\draw[gi interval] (2.7,0.2) -- (5.6,0.2) node[right] {$I_k$};
\draw[->] (0,-0.4) -- (6.6,-0.4);
\node[gi note] at (3.0,-0.7) {$r_1$};
\end{tikzpicture}
\end{graphillustration}


\subsection*{Графы сравнимости}
\begin{definition}[Граф сравнимости]
Пусть \(<\) — строгий частичный порядок на \(V\). Граф \(G=(V,\{uv: u<v\text{ или }v<u\})\)
называется \emph{графом сравнимости}. Эквивалентно: рёбра \(G\) можно ориентировать
\emph{транзитивно} (\(a\to b\to c\Rightarrow a\to c\)) и без циклов.
\end{definition}
Двудольные графы — графы сравнимости (все рёбра от одной доли к другой).
Индуцированный подграф графа сравнимости — граф сравнимости (ограничение порядка).
\begin{graphillustration}{Слева: частичный порядок (стрелка от меньшего к большему) и его граф сравнимости. Цвет $c(u)$ — число элементов наибольшей цепи, кончающейся в $u$ (цифра у вершины); цепь $a<c<e$ — клика размера $\omega=3$. Справа: «сеть» — не граф сравнимости.}{fig:13-comparability}
\begin{tikzpicture}[font=\small]
\begin{scope}
\node[gi vertex,label=left:{\scriptsize 1}] (a) at (0,0) {$a$};
\node[gi vertex,label=right:{\scriptsize 1}] (b) at (1.6,0) {$b$};
\node[gi vertex,label=left:{\scriptsize 2}] (c) at (0,1.3) {$c$};
\node[gi vertex,label=right:{\scriptsize 2}] (d) at (1.6,1.3) {$d$};
\node[gi vertex,label=left:{\scriptsize 3}] (e) at (0.8,2.6) {$e$};
\draw[gi arrow] (a) -- (c); \draw[gi arrow] (a) -- (d); \draw[gi arrow] (b) -- (d);
\draw[gi arrow] (c) -- (e); \draw[gi arrow] (d) -- (e);
\draw[gi arrow] (a) to[bend left=35] (e);
\draw[gi arrow] (b) to[bend right=35] (e);
\node[gi note] at (0.8,-1.2) {транзитивная ориентация};
\end{scope}
\begin{scope}[shift={(6,0.3)}]
\node[gi vertex] (a) at (0,0) {$a$};
\node[gi vertex] (b) at (1.8,0) {$c$};
\node[gi vertex] (c) at (0.9,1.5) {$b$};
\node[gi vertex] (a1) at (-1,-0.7) {$a'$};
\node[gi vertex] (b1) at (2.8,-0.7) {$c'$};
\node[gi vertex] (c1) at (0.9,2.7) {$b'$};
\draw[gi arrow] (a) -- (c); \draw[gi arrow] (c) -- (b); \draw[gi arrow] (a) -- (b);
\draw[gi edge] (a) -- (a1); \draw[gi edge] (b) -- (b1); \draw[gi edge] (c) -- (c1);
\node[gi note] at (0.9,-1.5) {ребро $bb'$ не ориентировать};
\end{scope}
\end{tikzpicture}
\end{graphillustration}

\begin{example}[Не граф сравнимости]
Треугольник \(abc\) с висячими вершинами \(a',b',c'\) («сеть»). Транзитивная ориентация
треугольника: \(a\to b\to c\), \(a\to c\) (с точностью до имён). Ребро \(bb'\): если \(b'\to b\),
то нужно \(b'\to c\); если \(b\to b'\), то нужно \(a\to b'\). Ни того, ни другого ребра нет.
\end{example}
\begin{theorem}[Мирский, 1971]
Граф сравнимости совершенен.
\end{theorem}
\begin{proof}
Индуцированные подграфы — тоже графы сравнимости, поэтому достаточно доказать
\(\chi(G)\le\omega(G)\). Клика графа сравнимости — \emph{цепь} (попарно сравнимые элементы),
так что \(\omega\) — длина наибольшей цепи. Положим
\[
c(u)=\max\bigl\{|L|:\ L\text{ — цепь, }u=\max L\bigr\}
\]
(в лекциях: длина наибольшего ориентированного пути, кончающегося в \(u\), плюс один).
Тогда \(c\colon V\to\{1,\ldots,\omega\}\). Если \(u<v\), любую цепь, оканчивающуюся в \(u\),
можно продолжить вершиной \(v\), поэтому \(c(v)\ge c(u)+1\). Значит, смежные вершины
получают разные цвета — правильная раскраска в \(\omega\) цветов.
\end{proof}
\begin{tcolorbox}[colback=cbGrayBack,colframe=cbGray!55]
\textbf{Дилуорс.} Дополнение графа сравнимости тоже совершенно — это следует из
слабой теоремы о совершенных графах (билет~\ref{ticket:14}). На языке порядков:
наибольшая антицепь равна наименьшему числу цепей, покрывающих множество.
\end{tcolorbox}

\ticket{Слабая теорема о совершенных графах}
\label{ticket:14}
\begin{theorem}[Ловас, 1972; слабая теорема о совершенных графах]
Граф \(G\) совершенен тогда и только тогда, когда \(\overline G\) совершенен.
\end{theorem}
Достаточно доказать «\(\Rightarrow\)» (обратное — то же для \(\overline G\), так как \(\overline{\overline G}=G\)).
Нам понадобятся три леммы.

\begin{lemma}[\(\alpha\chi\ge n\)]\label{lem:alphachi}
Для любого графа \(\alpha(G)\,\chi(G)\ge n\).
\end{lemma}
\begin{proof}
Покрасим \(G\) в \(\chi(G)\) цветов; классы цветов \(V_1,\ldots,V_{\chi}\) независимы, поэтому
\(n=\sum|V_i|\le\chi(G)\,\alpha(G)\).
\end{proof}

\begin{definition}[Клонирование вершины]
Для \(x\in V\) граф \(G\circ x\) получается добавлением новой вершины \(x'\), смежной с \(x\) и
со всеми соседями \(x\): \(G\circ x=\bigl(V\cup\{x'\},\;E\cup\{x'u: xu\in E\}\cup\{xx'\}\bigr)\).
\end{definition}
\begin{lemma}[Клонирование сохраняет совершенность]\label{lem:clone}
Если \(G\) совершенен, то \(G\circ x\) совершенен.
\end{lemma}
\begin{proof}
Индуцированный подграф \(G\circ x\) — либо индуцированный подграф \(G\) (если в нём нет
\(x\) или \(x'\); вершина \(x'\) без \(x\) играет роль \(x\)), либо \(H\circ x\) для индуцированного
\(H\le G\), \(x\in H\), который тоже совершенен. Поэтому достаточно доказать
\(\chi(G\circ x)\le\omega(G\circ x)\) для совершенного \(G\); пусть \(\omega=\omega(G)\).

\emph{1) \(x\) лежит в некоторой максимальной клике \(K\).} Тогда \(K\cup\{x'\}\) — клика, и
\(\omega(G\circ x)=\omega+1\). Покрасим \(G\) в \(\chi(G)=\omega\) цветов, а \(x'\) — в новый цвет:
\(\chi(G\circ x)\le\omega+1\).

\emph{2) \(x\) не лежит ни в одной максимальной клике.} Клика с \(x\) и \(x'\) имеет размер
\(|K|+1\le\omega\) для клики \(K\ni x\) графа \(G\), поэтому \(\omega(G\circ x)=\omega\). Покрасим \(G\) в
\(\omega\) цветов, пусть \(A\) — класс цвета \(x\). Каждая максимальная клика \(G\) содержит
вершины всех \(\omega\) цветов, значит, пересекает \(A\), причём не в \(x\) — т.\,е. пересекает
\(A\setminus\{x\}\). Поэтому \(\omega\bigl(G-(A\setminus\{x\})\bigr)\le\omega-1\), и по совершенности этот граф
красится в \(\omega-1\) цвет. Множество \((A\setminus\{x\})\cup\{x'\}\) независимо (\(x'\) смежна
только с \(x\) и \(N(x)\), а \(A\) — класс цвета \(x\)); дадим ему ещё один цвет. Итого \(\omega\) цветов.
\end{proof}
\begin{graphillustration}{Клонирование: $x'$ смежна с $x$ и со всеми соседями $x$ (штриховые рёбра). Справа случай 2: $A$ — класс цвета $x$; множество $(A\setminus\{x\})\cup\{x'\}$ независимо и получает отдельный цвет.}{fig:14-clone}
\begin{tikzpicture}[font=\small]
\begin{scope}
\node[gi vertex] (x) at (0,0) {$x$};
\node[gi vertex] (p) at (-1.2,1.2) {};
\node[gi vertex] (q) at (1.2,1.2) {};
\node[gi vertex] (r) at (0,-1.3) {};
\node[gi selected] (xx) at (1.4,-0.2) {$x'$};
\draw[gi edge] (x) -- (p); \draw[gi edge] (x) -- (q); \draw[gi edge] (x) -- (r); \draw[gi edge] (p) -- (q);
\draw[gi alternate] (xx) -- (x); \draw[gi alternate] (xx) -- (p); \draw[gi alternate] (xx) -- (q); \draw[gi alternate] (xx) -- (r);
\node[gi note] at (0,-2.0) {$G\circ x$};
\end{scope}
\begin{scope}[shift={(5.5,0)}]
\node[gi vertex] (x) at (0,0) {$x$};
\node[gi vertex] (a1) at (-1.6,-1.0) {};
\node[gi vertex] (a2) at (-0.5,-1.4) {};
\node[gi selected] (xx) at (1.4,-0.2) {$x'$};
\node[gi vertex] (n1) at (0,1.2) {};
\draw[gi edge] (x) -- (n1); \draw[gi alternate] (xx) -- (x); \draw[gi alternate] (xx) -- (n1);
\draw[gi edge] (a1) -- (n1); \draw[gi edge] (a2) -- (n1);
\node[gi group,fit=(x)(a1)(a2),label=left:{$A$}] {};
\node[gi note,anchor=west] at (2.2,0.2) {новый цвет:\\$(A\setminus\{x\})\cup\{x'\}$};
\end{scope}
\end{tikzpicture}
\end{graphillustration}


\begin{lemma}[Критерий через независимое множество]\label{lem:crit}
\(G\) совершенен \(\iff\) в каждом индуцированном \(H\le G\) есть независимое множество,
пересекающее все максимальные (размера \(\omega(H)\)) клики \(H\).
\end{lemma}
\begin{proof}
\(\Rightarrow\). Покрасим \(H\) в \(\chi(H)=\omega(H)\) цветов. Максимальная клика содержит вершины
всех цветов, поэтому класс любого цвета — искомое независимое множество.

\(\Leftarrow\). Докажем \(\chi(H)\le\omega(H)\) для всех \(H\le G\) индукцией по \(\omega(H)\) (лекционный
«алгоритм»: красим \(S\) в цвет \(i\), удаляем, \(i\gets i+1\)). При \(\omega(H)=0\) граф пуст. Пусть
\(S\) — независимое множество из условия. Каждая клика размера \(\omega(H)\) теряет вершину, поэтому
\(\omega(H-S)=\omega(H)-1\). Граф \(H-S\le G\) по индукции красится в \(\omega(H)-1\) цветов; \(S\) получает
ещё один.
\end{proof}

\begin{proof}[Доказательство теоремы Ловаса]
Пусть \(G\) совершенен. Клики \(\overline G\) — независимые множества \(G\), и наоборот;
индуцированные подграфы \(\overline G\) — дополнения индуцированных подграфов \(G\), а они
совершенны. По лемме~\ref{lem:crit}, применённой к \(\overline G\), достаточно доказать:
\begin{center}
\emph{в любом совершенном графе \(G\) есть клика \(K\), пересекающая все максимальные
независимые множества.}
\end{center}
Пусть \(\mathbb A\) — множество всех максимальных (размера \(\alpha=\alpha(G)\)) независимых
множеств, \(k=|\mathbb A|\), и для вершины \(v\)
\[
a_v=\bigl|\{I\in\mathbb A: v\in I\}\bigr| .
\]
Построим \(G'\): каждую вершину \(v\) заменим \(a_v\) попарно смежными копиями (клонируем
\(a_v-1\) раз; при \(a_v=0\) удаляем \(v\)). По лемме~\ref{lem:clone} и наследственности \(G'\) совершенен.

\emph{Число вершин.} \(v(G')=\sum_v a_v=\sum_v\sum_{I\in\mathbb A}[v\in I]=\sum_{I\in\mathbb A}|I|=k\alpha\).

\emph{\(\alpha(G')=\alpha\).} Независимое множество \(G'\) содержит не более одной копии каждой
вершины, и исходные вершины образуют независимое множество \(G\).

\emph{\(\chi(G')\le k\).} Распределим \(a_v\) копий \(v\) по одной между \(a_v\) множествами \(I\ni v\).
Копии, отнесённые к \(I\), образуют независимое множество (копия \(I\)); это раскраска
в \(k\) цветов, где цвет \(I\) состоит из копий вершин \(I\).

Из леммы~\ref{lem:alphachi}: \(k\alpha=v(G')\le\alpha(G')\chi(G')\le\alpha k\), значит,
\(\chi(G')=k\), и по совершенности \(\omega(G')=k\). Пусть \(K'\) — клика \(G'\) размера \(k\).
В раскраске выше \(k\) цветов, а вершины \(K'\) окрашены по-разному, поэтому \(K'\) содержит
ровно одну вершину каждого цвета.

Пусть \(K\) — множество вершин \(G\), копии которых лежат в \(K'\). Это клика (копии разных
вершин смежны, только если смежны оригиналы). Для любого \(I\in\mathbb A\) клика \(K'\)
содержит вершину цвета \(I\) — копию некоторой \(v\in I\); тогда \(v\in K\cap I\).
Итак, \(K\) пересекает все максимальные независимые множества.
\end{proof}
\begin{graphillustration}{Конструкция Ловаса для $G=P_4=abcd$. Максимальные независимые множества: $\{a,c\},\{a,d\},\{b,d\}$, $k=3$, $\alpha=2$; $a_a=a_d=2$, $a_b=a_c=1$. В $G'$ ровно $k\alpha=6$ вершин, клика $\{a_1,a_2,b\}$ размера $k$ проецируется в клику $K=\{a,b\}$, пересекающую все три множества.}{fig:14-blowup}
\begin{tikzpicture}[font=\small]
\begin{scope}
\node[gi vertex] (a) at (0,0) {$a$};
\node[gi vertex] (b) at (1.2,0) {$b$};
\node[gi vertex] (c) at (2.4,0) {$c$};
\node[gi vertex] (d) at (3.6,0) {$d$};
\draw[gi edge] (a) -- (b) -- (c) -- (d);
\node[gi note] at (1.8,-0.8) {$G=P_4$};
\end{scope}
\node[gi note] at (4.8,0) {$\Longrightarrow$};
\begin{scope}[shift={(6.2,0)}]
\node[gi selected] (a1) at (0,0.7) {$a_1$};
\node[gi selected] (a2) at (0,-0.7) {$a_2$};
\node[gi selected] (b) at (1.4,0) {$b$};
\node[gi vertex] (c) at (2.8,0) {$c$};
\node[gi vertex] (d1) at (4.2,0.7) {$d_1$};
\node[gi vertex] (d2) at (4.2,-0.7) {$d_2$};
\draw[gi match] (a1) -- (a2); \draw[gi match] (a1) -- (b); \draw[gi match] (a2) -- (b);
\draw[gi edge] (b) -- (c); \draw[gi edge] (c) -- (d1); \draw[gi edge] (c) -- (d2); \draw[gi edge] (d1) -- (d2);
\node[gi note] at (2.1,-1.4) {$G'$: $v(G')=6$, $\chi(G')=\omega(G')=3$};
\end{scope}
\end{tikzpicture}
\end{graphillustration}


\begin{tcolorbox}[colback=cbGrayBack,colframe=cbGray!55]
\textbf{Что ещё полезно знать (без доказательств).}
\begin{itemize}\tightlist
\item \emph{Критерий Ловаса:} \(G\) совершенен \(\iff\) \(\alpha(H)\,\omega(H)\ge v(H)\) для всех \(H\le G\);
  он симметричен относительно дополнения.
\item \emph{Сильная теорема (Чудновская–Робертсон–Сеймур–Томас, 2002; гипотеза Бержа):}
  \(G\) совершенен \(\iff\) ни \(G\), ни \(\overline G\) не содержат индуцированного нечётного цикла
  длины \(\ge5\) («нечётных дыр и антидыр»).
\item \emph{Хордальные графы} (нет индуцированных циклов длины \(\ge4\)) совершенны; это
  обобщает интервальные графы (билет~\ref{ticket:13}).
\item \emph{Следствие (Дилуорс):} дополнение графа сравнимости совершенно.
\end{itemize}
\end{tcolorbox}

\end{document}
